% Standalone source: compile this file directly with pdfLaTeX.
% No auxiliary project files are required.
\documentclass[11pt]{article}

\usepackage[a4paper,margin=1in]{geometry}
\usepackage[T1]{fontenc}
\usepackage[utf8]{inputenc}
\usepackage{lmodern}
\usepackage{microtype}
\usepackage{amsmath,amssymb,amsthm,mathtools}
\usepackage{mathrsfs}
\usepackage{xurl}
\usepackage{authblk}
\usepackage{xcolor}
\usepackage[
  colorlinks=true,
  linkcolor=blue!55!black,
  citecolor=blue!55!black,
  urlcolor=blue!55!black,
  unicode=true
]{hyperref}
\hypersetup{
  pdftitle={A Proposed Solution to Erd\H{o}s Problem 1002},
  pdfauthor={Shouqiao Wang}
}
\usepackage{enumitem}

\numberwithin{equation}{section}
\setlength{\emergencystretch}{2em}

\newtheorem{theorem}{Theorem}[section]
\newtheorem{proposition}[theorem]{Proposition}
\newtheorem{lemma}[theorem]{Lemma}
\newtheorem{corollary}[theorem]{Corollary}
\newtheorem{remark}[theorem]{Remark}
\theoremstyle{definition}
\newtheorem{definition}[theorem]{Definition}

\newcommand{\T}{\mathbb T}
\newcommand{\R}{\mathbb R}
\newcommand{\Z}{\mathbb Z}
\newcommand{\N}{\mathbb N}
\newcommand{\E}{\mathbb E}
\newcommand{\Pp}{\mathbb P}
\newcommand{\1}{\mathbf 1}
\newcommand{\e}{\mathrm e}
\newcommand{\ii}{\mathrm i}
\newcommand{\dd}{\,\mathrm d}
\newcommand{\Var}{\operatorname{Var}}
\newcommand{\TV}{\operatorname{TV}}
\newcommand{\PV}{\operatorname{PV}}
\newcommand{\nint}{\operatorname{nint}}
\newcommand{\sgn}{\operatorname{sgn}}
\newcommand{\Si}{\operatorname{Si}}
\newcommand{\lcm}{\operatorname{lcm}}
\newcommand{\Leb}{\operatorname{Leb}}
\newcommand{\norm}[1]{\left\lVert#1\right\rVert}
\newcommand{\abs}[1]{\left\lvert#1\right\rvert}
\newcommand{\braces}[1]{\left\{#1\right\}}
\newcommand{\angles}[1]{\left\langle#1\right\rangle}
\newcommand{\floor}[1]{\left\lfloor#1\right\rfloor}
\newcommand{\ceil}[1]{\left\lceil#1\right\rceil}
\newcommand{\wh}{\widehat}

\title{A Proposed Solution to Erd\H{o}s Problem 1002}
\author{Shouqiao Wang}
\affil{Columbia University\hspace{2.5em}Multiscalar Intelligence}
\date{}

\begin{document}
\maketitle

\begin{abstract}
Let
\[
 S_N(\alpha)=\sum_{k=1}^{N}\left(\frac12-\{k\alpha\}\right),
 \qquad \alpha\sim\operatorname{Unif}(0,1).
\]
We prove that
\[
 \frac{S_N(\alpha)}{\log N}\Longrightarrow
 \operatorname{Cauchy}\left(0,\frac1{2\pi}\right).
\]
Thus the distribution functions converge, at every real \(c\), to
\[
 \frac12+\frac1\pi\arctan(2\pi c).
\]
The starting point of the rotation is fixed; no averaging in a second
spatial coordinate or in time is used.  The proof first reconstructs the
sum in \(L^2\), up to \(o(\log N)\), as a primitive rational shot sum.
A Ramanujan-sum square-function argument removes all nonresonant shots.
The remaining signed, marked resonances converge by a continued-fraction
rare-event theorem proved here.  A cylinder oscillation argument supplies
the growing mark without averaging the starting point.  The resulting
Poisson integral has the stated Cauchy scale.
This proposed solution was found by GPT-5.6.
\end{abstract}

\section{Introduction}

Throughout, all logarithms are natural, and \(\Leb\) denotes Lebesgue
measure.  Erd\H{o}s asked \cite[p.~63]{Erdos1964} whether the random variables
\[
 S_N(\alpha):=\sum_{k=1}^{N}\left(\frac12-\{k\alpha\}\right),
 \qquad f_N(\alpha):=\frac{S_N(\alpha)}{\log N},
 \qquad \alpha\sim\operatorname{Unif}(0,1),
\]
possess an asymptotic distribution function along the full sequence of
integers \(N\), with convergence at every real threshold.  We give an
affirmative answer and identify the limit explicitly as a centered Cauchy
law.

\begin{theorem}[Main theorem]\label{thm:main}
For every \(c\in\R\),
\[
 \lim_{N\to\infty}
 \Leb\left\{\alpha\in(0,1):
 \frac{S_N(\alpha)}{\log N}\le c\right\}
 =
 g(c):=\frac12+\frac1\pi\arctan(2\pi c).
\]
Equivalently,
\[
 \frac{S_N}{\log N}\Longrightarrow X,\qquad
 \E e^{\ii tX}=\exp\left(-\frac{|t|}{2\pi}\right).
\]
\end{theorem}

Write
\[
 e(x)=\exp(2\pi\ii x),\qquad
 \psi(x)=\frac12-\{x\},\qquad
 V(x)=\frac{\{x\}(1-\{x\})}{2}.
\]
We write
\[
 \norm{u}_{\T}:=\min_{j\in\Z}|u-j|,\qquad
 \Si(x):=\int_0^x\frac{\sin t}{t}\dd t,
\]
with the integrand at zero defined by continuity.  The symbol
\(\nint(t)\) denotes a nearest integer to \(t\), with the smaller one
chosen at a tie.  Values at the finitely many rational discontinuities
will never affect an integral or a distribution.  We identify
\(\T=\R/\Z\) with \([0,1)\).

The distinction between Theorem~\ref{thm:main} and the classical
two-variable result is important.  Kesten's spatial theorem averages both
the rotation parameter and the starting point
\cite{Kesten1960,Kesten1962}; here the starting point is exactly zero.

For \(p\ge1\), let
\[
 q_p=q_p(\alpha):=\nint(p\alpha),\qquad
 \delta_p=\delta_p(\alpha):=p\alpha-q_p\in(-1/2,1/2],
\]
put \(L=L_N:=\log N\), and impose \((p,q_p)=1\).  The central object is
\begin{equation}\label{eq:defY}
 Y_N(\alpha)=
 \sum_{\substack{1\le p\le N\\(p,q_p)=1}}
 \frac{V(N\delta_p)}{p\delta_p}.
\end{equation}
At the null set where \(\delta_p=0\), this quotient is assigned the value
zero.  This convention will be used for every shot sum below.
The proof has three stages.  The final transition-band step in
stage~(ii) invokes the marked Poisson theorem from stage~(iii); that
theorem is proved independently of stage~(ii), so there is no circular
dependence.

\begin{enumerate}[label=(\roman*)]
\item In Section~\ref{sec:reconstruction} we prove
\begin{equation}\label{eq:outline-reconstruction}
 \norm{S_N-Y_N}_2=o(\log N).
\end{equation}
The proof is an exact Fourier--Ramanujan reconstruction, including its
principal-value and endpoint conventions.

\item In Section~\ref{sec:minor} we prove the small-jump estimate: for
every \(\eta>0\),
\begin{equation}\label{eq:outline-minor}
 \lim_{A\to\infty}\limsup_{N\to\infty}
 \Leb\left\{
 \frac1{\log N}\left|
 \sum_{\substack{p\le N,\ (p,q_p)=1\\
       (\log N)|p\delta_p|>A}}
 \frac{V(N\delta_p)}{p\delta_p}\right|>\eta\right\}=0.
\end{equation}
This is the fixed-start step which is absent from the usual
two-variable argument.

\item In Section~\ref{sec:poisson}, the complementary finite process
\[
 \left\{\big((\log N)p\delta_p,\;N\delta_p\bmod 1\big):
 p\le N,\ (p,q_p)=1\right\}
\]
is identified with rare regular continued-fraction convergents.  Exponential
mixing proves the unmarked Poisson law; a two-scale cylinder argument proves
that the marks become independent Haar variables.  The limiting marked
process has intensity \(1/\zeta(2)\), and sending the compact cutoff to infinity gives
the characteristic exponent \(-|t|/(2\pi)\).
\end{enumerate}

We use the Fourier convention
\[
 \wh f(n)=\int_0^1 f(\alpha)e(-n\alpha)\dd\alpha.
\]
For an integrable function on \(\R\), the corresponding real-line
convention is
\[
 \wh f(t)=\int_{\R}f(v)e(-tv)\dd v.
\]
All \(L^2\)-norms without a displayed measure are over \(\T\).
We use the standard arithmetic functions \(\mu\) (Möbius), \(\varphi\)
(Euler), \(\tau\) (the divisor-counting function), and
\(\sigma_s(n)=\sum_{d\mid n}d^s\).

\section{Primitive-shot reconstruction}\label{sec:reconstruction}

\subsection{Periodization and its Fourier transform}

The absolutely convergent Fourier series
\begin{equation}\label{eq:Vfourier}
 V(x)=\frac1{12}-\sum_{\ell\ne0}
 \frac{e(\ell x)}{4\pi^2\ell^2}
 =\sum_{r\ge1}\frac{1-\cos(2\pi r x)}{2\pi^2r^2}
\end{equation}
will provide the mark.  For \(P\ge1\), define the nearest-cell sum
\[
 Y_{N,P}(\alpha)=
 \sum_{\substack{p\le P\\(p,q_p)=1}}
 \frac{V(N\delta_p)}{p\delta_p}
\]
and its full periodization
\begin{equation}\label{eq:defZ}
 Z_{N,P}(\alpha)=
 \sum_{p\le P}\frac1p\,
 \PV\sum_{\substack{q\in\Z\\(p,q)=1}}
 \frac{V(N(p\alpha-q))}{p\alpha-q}.
\end{equation}
The symmetric principal value is taken in \(q\).

\begin{lemma}[Fourier transform]\label{lem:transform}
For \(s>0\),
\begin{equation}\label{eq:Iinfty}
 I_N(s):=\PV\int_{\R}\frac{V(Nx)}x e(-sx)\dd x
 =-\frac{\ii}{2\pi}\sum_{k>s/N}\frac1{k^2},
\end{equation}
where the term \(k=s/N\) has weight \(1/2\) when \(s/N\) is an
integer.  Consequently, for \(n>0\),
\begin{equation}\label{eq:Zcoeff}
 \wh Z_{N,P}(n)=
 -\frac{\ii}{2\pi}
 \sum_{p\le P}\frac{c_p(n)}{p^2}
 H_*\left(\frac{n}{pN}\right),
 \qquad
 H_*(u)=\sum_{k>u}\frac1{k^2}
 +\frac{\1_{\{u\in\N\}}}{2u^2},
\end{equation}
where the last term is understood as zero at \(u=0\), and
\[
 c_p(n)=\sum_{\substack{a\bmod p\\(a,p)=1}}e(an/p)
\]
is the Ramanujan sum.
Moreover,
\[
 \widehat Z_{N,P}(0)=0,\qquad
 \widehat Z_{N,P}(-n)=\overline{\widehat Z_{N,P}(n)}.
\]
\end{lemma}

\begin{proof}
Put \(h_N(x)=V(Nx)/x\), with an arbitrary value at \(x=0\).
Before computing, we make the periodization limit explicit.  Let
\(\mathcal A_p\) be the reduced residues in \(\{0,\ldots,p-1\}\), and
for \(R\ge1\) define the finite sum
\[
 Z_{N,P}^{(R)}(\alpha)=
 \sum_{p\le P}\frac1p
 \sum_{a\in\mathcal A_p}\sum_{|j|\le R}
 h_N(p\alpha-a-pj).
\]
For \(z=p\alpha-a\in[-p,p]\), periodicity of \(V\) gives, when
\(j\ge2\),
\[
 h_N(z-pj)+h_N(z+pj)
 =
 V(Nz)\left(\frac1{z-pj}+\frac1{z+pj}\right),
\]
and hence
\[
 \sup_{|z|\le p}
 \left|\sum_{j>R}\bigl(h_N(z-pj)+h_N(z+pj)\bigr)\right|
 \ll \sum_{j>R}\frac{|z|}{p^2j^2}
 \ll\frac1R.
\]
The finitely many apparent quotient singularities are canceled by the
zeros of the numerator and leave only bounded one-sided jumps, because
\(V(Nx)=O_N(|x|)\) at an integer.  Thus
\(Z_{N,P}^{(R)}\) converges in \(L^2(\T)\), in fact uniformly away
from those finitely many jump points, to \(Z_{N,P}\).  Compare this
residue-wise cutoff \(|j|\le R\) with the literal symmetric cutoff
\(|q|\le pR\).  In a fixed residue class \(q=a+pj\), the two index
intervals differ by at most two endpoint indices, and the corresponding
summands are \(O_p(R^{-1})\).  Thus their difference is \(O_p(R^{-1})\).
For an arbitrary literal cutoff radius \(S\), take
\(R=\lfloor S/p\rfloor\).  The cutoffs \(|q|\le S\) and
\(|q|\le pR\) differ by \(O_p(1)\) endpoints, each of size
\(O_p(S^{-1})\).  Hence the residue-wise convention and the symmetric
principal value in \(q\) used in \eqref{eq:defZ} have the same limit
(with \(p\le P\) fixed here).

To justify the principal-value calculation, first integrate over
\(\eta<|x|<R\), insert the absolutely and uniformly convergent series
\eqref{eq:Vfourier}, and then let \(\eta\downarrow0\) symmetrically and
\(R\to\infty\).  The first limit can be passed through the series because,
uniformly near zero,
\[
 \sum_{k\ge1}\frac{|1-\cos(2\pi kNx)|}{k^2|x|}\ll N;
\]
split the sum at \(k=(N|x|)^{-1}\).  For the second limit, the symmetric
truncated integrals of each exponential divided by \(x\) are bounded
uniformly in its frequency and in \(R\), by the elementary boundedness of
the sine integral.  Thus \(\sum k^{-2}<\infty\) supplies domination.
Equivalently one may insert \(e^{-\epsilon|x|}\), retain the symmetric
principal value at zero, and let \(\epsilon\downarrow0\); for every real
\(a\),
\[
 \lim_{\epsilon\downarrow0}
 \PV\int_{\R}\frac{e(-ax)e^{-\epsilon|x|}}x\dd x
 =-\ii\pi\sgn(a),
\]
with value zero at \(a=0\).
The distributional identity
\[
 \PV\int_{\R}\frac{e(-sx)}x\dd x=-\ii\pi\sgn(s)
\]
and the second series in \eqref{eq:Vfourier} show that the \(k\)-th
summand contributes \(-\ii/(2\pi k^2)\) precisely when \(kN>s\),
zero when \(kN<s\), and half at equality.  This proves
\eqref{eq:Iinfty}.

We may now compute with the finite periodization.  The changes of
variables \(t=p\alpha\) and \(x=t-a-pj\) give, for every integer \(n\ge0\),
\[
 \widehat Z_{N,P}^{(R)}(n)
 =
 \sum_{p\le P}\frac1{p^2}
 \sum_{a\in\mathcal A_p}e(-na/p)
 \int_{-a-pR}^{\,p-a+pR}h_N(x)e(-nx/p)\dd x.
\]
Indeed, as \(j\) ranges from \(-R\) to \(R\), the corresponding
intervals of length \(p\) are adjacent and tile the displayed interval.
Replacing its endpoints by \(\pm pR\) changes the integral by
\(O_p(R^{-1})\), since \(|h_N(x)|\ll |x|^{-1}\) there.  Hence, for
\(n>0\), the integral tends to \(I_N(n/p)\); for \(n=0\), it tends
to zero because \(h_N\) is odd and the same endpoint estimate reduces
the limit to a symmetric integral.  Since
\[
 \sum_{a\in\mathcal A_p}e(-na/p)=c_p(-n)=c_p(n),
 \]
the \(L^2\)-limit \(Z_{N,P}^{(R)}\to Z_{N,P}\) and continuity of
Fourier coefficients on \(L^2\) prove \eqref{eq:Zcoeff} and
\(\widehat Z_{N,P}(0)=0\).  Finally \(Z_{N,P}\) is real, which gives
the negative coefficients by conjugacy.
\end{proof}


\subsection{The exact all-denominator identity}

Define \(Z_{N,\infty}\) to be the mean-zero \(L^2(\T)\) function whose
positive Fourier coefficients are given by the right side of
\eqref{eq:Zcoeff} with the sum over all \(p\ge1\), and whose negative
coefficients are their conjugates.  The series over \(p\) is absolutely
convergent for each coefficient, since the divisor formula for Ramanujan
sums gives
\[
 \sum_{p\ge1}\frac{|c_p(n)|}{p^2}
 \le \sum_{d\mid n}\frac1d\sum_{r\ge1}\frac1{r^2}
 \ll \sigma_{-1}(n).
\]
Lemma~\ref{lem:collapse} below shows that
the resulting positive coefficient is
\[
 O_N(n^{-1}),
\]
and hence that the coefficient sequence is square summable.  This also fixes the
symmetric principal-value convention in the notation \(Z_{N,\infty}\).

\begin{lemma}[Möbius collapse]\label{lem:collapse}
For every \(n>0\),
\begin{equation}\label{eq:collapse}
 \sum_{p\ge1}\frac{c_p(n)}{p^2}
 H_*\left(\frac{n}{pN}\right)
 =
 \frac1n\left(
 \sum_{\substack{d\mid n\\d<N}}d+
 \frac N2\1_{\{N\mid n\}}\right).
\end{equation}
Hence, in \(L^2(\T)\),
\begin{equation}\label{eq:exactZ}
 Z_{N,\infty}=S_N-\frac12\psi(N\alpha).
\end{equation}
\end{lemma}

\begin{proof}
Use
\[
 c_p(n)=\sum_{a\mid(p,n)}a\,\mu(p/a)
\]
and write \(p=ar\).  The left side of \eqref{eq:collapse} becomes
\[
 \sum_{a\mid n}\frac1a
 \sum_{\substack{r,k\ge1\\rk>n/(aN)}}
 \frac{\mu(r)}{r^2k^2}.
\]
On grouping by \(m=rk\),
\[
 \sum_{rk>X}\frac{\mu(r)}{r^2k^2}
 =\sum_{m>X}\frac1{m^2}\sum_{r\mid m}\mu(r)
 =
 \begin{cases}
 1,&X<1,\\
 0,&X>1,
 \end{cases}
\]
with midpoint value \(1/2\) at \(X=1\).  Replacing \(a\) by
\(d=n/a\) gives \eqref{eq:collapse}.

Since
\[
 \psi(x)=\sum_{m\ne0}\frac{e(mx)}{2\pi\ii m}
\quad\text{in }L^2(\T),
\]
the positive \(n\)-th coefficient of \(S_N\) is
\[
 -\frac{\ii}{2\pi n}\sum_{\substack{d\mid n\\d\le N}}d.
\]
Combining this with Lemma~\ref{lem:transform} and
\eqref{eq:collapse} yields exactly the coefficient of
\(S_N-\frac12\psi(N\alpha)\).  Negative coefficients are conjugate,
and both sides have mean zero.
\end{proof}

\subsection{Nearest cell versus full periodization}

\begin{proposition}[Window estimate]\label{prop:window}
Uniformly for \(N,P\ge2\),
\begin{equation}\label{eq:window}
 \norm{Y_{N,P}-Z_{N,P}}_2^2
 \ll \log(2P)\bigl(1+\log\log(e^e+NP)\bigr).
\end{equation}
\end{proposition}

\begin{proof}
Let \(q(t)=\nint(t)\).  Möbius inversion and the cotangent
principal-value identity give
\begin{equation}\label{eq:mobcot}
 \PV\sum_{\substack{q\in\Z\\(p,q)=1}}\frac1{t-q}
 =\sum_{d\mid p}\mu(d)\frac{\pi}{d}\cot\left(\frac{\pi t}{d}\right).
\end{equation}
Put
\[
 L_d(u)=
 \frac{\1_{\{\norm{u}_{\T}<1/(2d)\}}}{u-\nint(u)}
 -\pi\cot(\pi u),
\]
where the first quotient is interpreted in the unique local lift.  The
poles cancel.  Since \(V(N(t-q))=V(Nt)\), writing \(p=dr\) in the
difference between the nearest coprime pole and \eqref{eq:mobcot} yields
\begin{equation}\label{eq:window-decomp}
 Y_{N,P}-Z_{N,P}
 =\sum_{dr\le P}\frac{\mu(d)}{d^2r}F_{N,d}(r\alpha),
 \qquad
 F_{N,d}(u)=V(Ndu)L_d(u).
\end{equation}

For \(m\ne0\), direct integration on
\((-1/(2d),1/(2d))\) gives
\begin{equation}\label{eq:Lhat}
 \wh L_d(m)=
 \ii\pi\sgn(m)-2\ii\Si(\pi m/d),\qquad
 \abs{\wh L_d(m)}\ll\min\left(1,\frac d{|m|}\right),
\end{equation}
and \(\wh L_d(0)=0\).  If \(v_\ell\) are the coefficients in
\eqref{eq:Vfourier}, then
\[
 \wh F_{N,d}(m)=\sum_{\ell\in\Z}v_\ell
 \wh L_d(m-\ell Nd).
\]
For fixed \(\ell\), the contribution to the \(n\)-th Fourier coefficient
of \eqref{eq:window-decomp} is
\begin{equation}\label{eq:C-ell}
 C_\ell(n)=
 \sum_{\substack{dr\le P\\r\mid n}}
 \frac{\mu(d)}{d^2r}
 \wh L_d(n/r-\ell Nd).
\end{equation}

We use the elementary estimate
\begin{equation}\label{eq:Lweighted}
 \sum_{m\ge1}\sigma_{-1}(m)
 \abs{\wh L_d(m-M)}^2
 \ll
 d\bigl(1+\log\log(e^e+|M|+d)\bigr).
\end{equation}
Indeed, on \(|m-M|\le d\), the first bound in
\eqref{eq:Lhat} and
\(\sigma_{-1}(m)\ll1+\log\log(e^e+m)\) give
\[
 \sum_{\substack{m\ge1\\|m-M|\le d}}
 \sigma_{-1}(m)|\widehat L_d(m-M)|^2
 \ll d\bigl(1+\log\log(e^e+|M|+d)\bigr).
\]
On the annulus \(2^jd<|m-M|\le2^{j+1}d\), the second bound in
\eqref{eq:Lhat} gives
\[
 \sum_{\substack{m\ge1\\2^jd<|m-M|\le2^{j+1}d}}
 \sigma_{-1}(m)|\widehat L_d(m-M)|^2
 \ll 2^{-j}d
 \bigl(1+\log\log(e^e+|M|+2^{j+1}d)\bigr).
\]
Here
\[
 \log\log(e^e+|M|+2^{j+1}d)
 \ll \log\log(e^e+|M|+d)+\log(j+2),
\]
and both \(\sum_{j\ge0}2^{-j}\) and
\(\sum_{j\ge0}2^{-j}\log(j+2)\) converge.  Summing the annuli proves
\eqref{eq:Lweighted}.

Apply Cauchy--Schwarz to \eqref{eq:C-ell} with the fixed weights
\(d^{-3/2}r^{-1}\).  Since
\[
 \sum_{\substack{dr\le P\\r\mid n}}d^{-3/2}r^{-1}
 \ll\sigma_{-1}(n),
\]
this gives
\[
 |C_\ell(n)|^2
 \ll \sigma_{-1}(n)
 \sum_{\substack{dr\le P\\r\mid n}}
 d^{-5/2}r^{-1}
 \left|\widehat L_d(n/r-\ell Nd)\right|^2.
\]
Sum over \(n=rm\), use
\(\sigma_{-1}(rm)\le\sigma_{-1}(r)\sigma_{-1}(m)\), and then
\eqref{eq:Lweighted} with \(M=\ell Nd\).  We obtain
\[
 \sum_{n\ge1}|C_\ell(n)|^2
 \ll
 \log(2P)\bigl(1+\log\log(e^e+|\ell|NP)\bigr),
\]
because
\(\sum_{r\le R}\sigma_{-1}(r)/r\ll\log(2R)\) and
\(\sum_d d^{-3/2}<\infty\).
The negative frequencies are conjugate, while the zero coefficient
vanishes because \(V(Ndu)\) is even and \(L_d(u)\) is odd.
Finally,
\[
 \sum_{\ell\in\Z}|v_\ell|
 \sqrt{1+\log\log(e^e+|\ell|NP)}
 \ll
 \sqrt{1+\log\log(e^e+NP)}
\]
by \(|v_\ell|\ll(1+\ell^2)^{-1}\).  Minkowski and Plancherel
prove \eqref{eq:window}.
\end{proof}

\subsection{The superlinear tail and the natural cutoff}

\begin{lemma}[Crude all-\(p\) tail]\label{lem:allp-tail}
For \(P\ge N\),
\begin{equation}\label{eq:allp-tail}
 \norm{Z_{N,\infty}-Z_{N,P}}_2^2
 \ll \frac NP\bigl(1+\log(NP)\bigr)^3.
\end{equation}
\end{lemma}

\begin{proof}
Expanding \(c_p(n)\), the positive \(n\)-th coefficient, apart from
the fixed factor \(-\ii/(2\pi)\), is
\[
 A_P(n)=
 \sum_{d\mid n}\frac1d
 \sum_{r>P/d}\frac{\mu(r)}{r^2}
 H_*\left(\frac{n}{drN}\right).
\]
Since \(H_*(u)\ll\min(1,u^{-1})\), dropping the Möbius signs gives
\[
 |A_P(n)|\ll
 \begin{cases}
 \tau(n)/P,&n\le NP,\\[2mm]
 \displaystyle
 \frac{N\tau(n)}n\left(1+\log\frac{n}{NP}\right),&n>NP.
 \end{cases}
\]
The first line follows by summing \(r^{-2}\) above \(P/d\).
For the second, split the \(r\)-sum at \(n/(dN)\).  Now use
\[
 \sum_{n\le X}\tau(n)^2\ll X(1+\log X)^3
\]
to obtain
\[
 \sum_{n\le NP}|A_P(n)|^2
 \ll \frac1{P^2}\sum_{n\le NP}\tau(n)^2
 \ll \frac NP\bigl(1+\log(NP)\bigr)^3.
\]
For \(2^jNP<n\le2^{j+1}NP\), \(j\ge0\), the second pointwise
bound similarly gives
\[
 \sum_{2^jNP<n\le2^{j+1}NP}|A_P(n)|^2
 \ll \frac NP\,2^{-j}(j+2)^2
       \bigl(1+\log(2^{j+1}NP)\bigr)^3.
\]
The last expression is summable in \(j\), with total
\(O((N/P)(1+\log(NP))^3)\).  Summing the dyadic blocks proves
\eqref{eq:allp-tail}.
\end{proof}

We next improve the denominator cutoff from
\(N(1+\log N)^{10}\) to \(N\).  The following elementary
incomplete orthogonality will be used twice.

\begin{lemma}[Incomplete Ramanujan orthogonality]\label{lem:incomplete}
If \(p\ne p'\), then for every finite interval \(I\subset\Z\),
\begin{equation}\label{eq:incomplete}
 \left|\sum_{n\in I}c_p(n)c_{p'}(n)\right|
 \le2\sigma_1(p)\sigma_1(p').
\end{equation}
\end{lemma}

\begin{proof}
Insert
\[
 c_p(n)=\sum_{d\mid(p,n)}d\,\mu(p/d)
\]
twice.  The number of multiples of \(\lcm(d,e)\) in \(I\) is
\(|I|/\lcm(d,e)+O(1)\).  The main term vanishes when
\(p\ne p'\), by complete-period orthogonality of Ramanujan sums.  Indeed,
averaging their exponential definitions over one period
\(\lcm(p,p')\) leaves only pairs of reduced fractions satisfying
\(a/p+b/p'\in\Z\); reduction forces \(p=p'\).
The absolute sum of the endpoint errors is at most the right side of
\eqref{eq:incomplete}.
\end{proof}

\begin{lemma}[Discrete Abel summation against a BV weight]
\label{lem:discrete-BV}
Let \((a_n)\) be indexed by \(\Z\), or by a one-sided integer interval,
and suppose that
\begin{equation}\label{eq:discrete-BV-hypothesis}
 \sup_I\left|\sum_{n\in I}a_n\right|\le M,
\end{equation}
where the supremum is over all finite integer intervals in the index
set.  If \(W\) is a regulated function of bounded variation and tends
to zero at every infinite end of the index set, then
\begin{equation}\label{eq:discrete-BV}
 \left|\sum_n a_nW(n)\right|
 \le 2M\bigl(\norm{W}_\infty+\TV_{\R}(W)\bigr).
\end{equation}
The same conclusion holds with \(\TV_{\R}(W)\) replaced by the
discrete variation \(\sum_n|W(n+1)-W(n)|\).
\end{lemma}

\begin{proof}
For integers \(u\le v\), put
\(A_u(k)=\sum_{u\le n\le k}a_n\).  Ordinary summation by parts gives
the exact identity
\begin{equation}\label{eq:discrete-Abel-identity}
 \sum_{n=u}^v a_nW(n)
 =A_u(v)W(v)+
  \sum_{n=u}^{v-1}A_u(n)\bigl(W(n)-W(n+1)\bigr).
\end{equation}
By \eqref{eq:discrete-BV-hypothesis}, its absolute value is at most
\[
 M\left(|W(v)|+
 \sum_{n=u}^{v-1}|W(n+1)-W(n)|\right).
\]
On either infinite tail, both the endpoint value and the remaining
variation tend to zero.  Thus the finite sums satisfy the Cauchy
criterion.  Letting the endpoints tend to the ends of the index set in
\eqref{eq:discrete-Abel-identity} proves \eqref{eq:discrete-BV}; the
displayed factor \(2\) harmlessly covers all one-sided and two-sided
endpoint conventions.
\end{proof}

\begin{proposition}[Natural denominator cutoff]\label{prop:natural-cutoff}
As \(N\to\infty\),
\begin{equation}\label{eq:natural-cutoff-Z}
 \norm{Z_{N,\infty}-Z_{N,N}}_2=o(\log N),
\end{equation}
and consequently
\begin{equation}\label{eq:S-Y}
 \boxed{\norm{S_N-Y_N}_2=o(\log N).}
\end{equation}
\end{proposition}

\begin{proof}
Put \(L=\log N\) and \(P=N(1+L)^{10}\).  For a dyadic block
\(Q<p\le2Q\), \(Q>N\), let
\[
 A_Q(n)=\sum_{Q<p\le2Q}\frac{c_p(-n)}{p^2}
 H_*\left(\frac{n}{pN}\right),\qquad n>0.
\]
Write \(h_p(n)=H_*(n/(pN))\).  Expanding the square gives
\begin{equation}\label{eq:natural-block-expansion}
 \sum_{n\ge1}|A_Q(n)|^2
 =\sum_{p,p'\asymp Q}\frac1{p^2p'^2}
  \sum_{n\ge1}c_p(-n)c_{p'}(-n)h_p(n)h_{p'}(n).
\end{equation}
The midpoint convention in \(H_*\) preserves monotonicity.  More
precisely, at every integer \(m\ge1\),
\[
 H_*(m^-)-H_*(m)=H_*(m)-H_*(m^+)=\frac1{2m^2}.
\]
Consequently each \(h_p\) is nonnegative and nonincreasing, tends to
zero, and is bounded by \(\zeta(2)\).  Hence \(h_ph_{p'}\) is also
nonnegative and nonincreasing, and
\[
 \norm{h_ph_{p'}}_\infty+
 \sum_{n\ge1}|h_p(n+1)h_{p'}(n+1)-h_p(n)h_{p'}(n)|
 \le2\zeta(2)^2.
\]
For \(p\ne p'\), apply Lemma~\ref{lem:discrete-BV} to
\(a_n=c_p(-n)c_{p'}(-n)\), and use
Lemma~\ref{lem:incomplete}.  Thus the inner sum in
\eqref{eq:natural-block-expansion} is
\(O(\sigma_1(p)\sigma_1(p'))\), including every midpoint jump and
every Abel boundary term.  The off-diagonal contribution is therefore
\[
 \sum_{\substack{p,p'\asymp Q\\p\ne p'}}
 \frac{\sigma_1(p)\sigma_1(p')}{p^2p'^2}
 \ll(\log\log(3Q))^2.
\]
For the diagonal, use
\[
 \sum_{n\bmod p}c_p(n)^2=p\varphi(p),\qquad
 \sum_{n\ge1}H_*(n/(pN))^2\ll pN.
\]
More explicitly, on the \(j\)-th period \(jp<n\le(j+1)p\), monotonicity
of \(H_*\) and complete-period orthogonality give a contribution at most
\[
 p\varphi(p)H_*(j/N)^2.
\]
Since \(H_*(u)\ll\min(1,u^{-1})\),
\(\sum_{j\ge0}H_*(j/N)^2\ll N\).  Thus subdivision into periods of
length \(p\) gives
\[
 \sum_{p\asymp Q}\frac1{p^4}
 \sum_{n\ge1}c_p(n)^2H_*(n/(pN))^2
 \ll\sum_{p\asymp Q}\frac{N\varphi(p)}{p^3}
 \ll\frac NQ.
\]
Thus
\begin{equation}\label{eq:high-block}
 \norm{A_Q}_{\ell^2(n\ge1)}^2
 \ll1+(\log\log(3Q))^2.
\end{equation}
There are \(O(\log(P/N))=O(\log L)\) blocks.  Minkowski,
\eqref{eq:high-block}, and Lemma~\ref{lem:allp-tail} give
\[
 \norm{Z_{N,\infty}-Z_{N,N}}_2
 \ll(\log L)^2+o(1)=o(L),
\]
which is \eqref{eq:natural-cutoff-Z}.

Finally combine \eqref{eq:exactZ}, Proposition~\ref{prop:window} with
\(P=N\), and \eqref{eq:natural-cutoff-Z}.  The term
\(\frac12\psi(N\alpha)\) has bounded \(L^2\)-norm, while
\[
 \norm{Y_{N,N}-Z_{N,N}}_2
 \ll\sqrt{L\log L}=o(L).
\]
This proves \eqref{eq:S-Y}.
\end{proof}

\section{The nonresonant shots}\label{sec:minor}

Throughout this section \(L=\log N\), and \(n\asymp K\) always means
\(K<n\le2K\), with \(K\) ranging over dyadic numbers.  We first control the range
\[
 \frac A L<|p\delta_p|<\varepsilon
\]
with a smooth cutoff, and then the fixed-away range
\(|p\delta_p|\ge\varepsilon\).  The two arguments are different.

\subsection{A mean-square Ramanujan tail}

We use the following consequence of Chan and Kumchev's mean-square
theorem for partial sums of Ramanujan sums
\cite[Theorem 1.2]{ChanKumchev}.

\begin{lemma}[Ramanujan tail]\label{lem:ram-tail}
There is an absolute \(B_{\rm tail}>30\) such that, uniformly for \(K\ge3\)
and \(2\le Q\le \sqrt K\),
\begin{equation}\label{eq:ram-tail}
 \left\|
 \sum_{q>Q}\frac{c_q(n)}{q^2}
 \right\|_{\ell^2(1\le n\le2K)}
 \ll\frac{\sqrt K}{Q}.
\end{equation}
The same bound holds for a finite tail
\(\sum_{Q<q\le R}\), uniformly in \(R\).
\end{lemma}

\begin{proof}
Put \(C_x(n)=\sum_{q\le x}c_q(n)\).
Set
\[
 x_{\rm CK}=x,\qquad y_{\rm CK}=2K,\qquad B_{\rm CK}=11.
\]
When \(x\le\sqrt {y_{\rm CK}}\), their elementary formula
\[
 \sum_{n\le y_{\rm CK}}|C_x(n)|^2
 =\frac{y_{\rm CK}x^2}{2\zeta(2)}
  +O(x^4+xy_{\rm CK}\log x)
\]
is \(O(Kx^2)\).  When
\(\sqrt {y_{\rm CK}}<x\le T:=K/(\log K)^{B_{\rm tail}}\), apply
part~(ii) of their theorem with the displayed substitution.  We verify
both of its range hypotheses, including the lower logarithmic margin.
Since \(B_{\rm CK}=11\), \(B_{\rm tail}>30\), and \(x\le T\),
\[
 x(\log x)^{2B_{\rm CK}}
 \le T(\log T)^{22}
 \le K(\log K)^{22-B_{\rm tail}}
 \le 2K=y_{\rm CK}
\]
for all sufficiently large \(K\).  Since
\(x>\sqrt {y_{\rm CK}}\), we also have
\[
 y_{\rm CK}<x^2\le x^2(\log x)^{B_{\rm CK}}.
\]
The bounded remaining values of \(K\) are absorbed in the constant.
Moreover, \(\log(xy_{\rm CK})\asymp\log x\asymp\log K\) throughout
this range.  The function denoted by \(\kappa\) in that theorem
satisfies \(\sup_{u\in\R}|\kappa(u)|<\infty\), so its main term is
\(O(y_{\rm CK}x^2)=O(Kx^2)\).  Its error relative to
\(y_{\rm CK}x^2\) is
\[
 O\!\left((\log x)^{10}
 \left(x^{-1/2}+(y_{\rm CK}/x)^{-1/2}\right)\right)=O(1)
\]
when \(B_{\rm tail}>30\): indeed
\(x^{-1/2}\le y_{\rm CK}^{-1/4}\), and
\(y_{\rm CK}/x\ge2(\log K)^{B_{\rm tail}}\) at the upper endpoint of the
range.  Thus
\begin{equation}\label{eq:CK}
 \sum_{n\le2K}|C_x(n)|^2\ll Kx^2
 \qquad
 (2\le x\le K/(\log K)^{B_{\rm tail}}).
\end{equation}
Partial summation gives
\[
 \sum_{Q<q\le R}\frac{c_q(n)}{q^2}
 =
 \frac{C_R(n)}{R^2}-\frac{C_Q(n)}{Q^2}
 +2\int_Q^R\frac{C_x(n)}{x^3}\dd x.
\]
Minkowski and \eqref{eq:CK} bound the part below
\(T=K/(\log K)^{B_{\rm tail}}\) by \(C\sqrt K/Q\).
Above \(T\), the divisor formula gives
\[
 |C_x(n)|\le x\tau(n),\qquad
 \norm{C_x}_{\ell^2(n\le2K)}
 \ll x\sqrt K(\log K)^{3/2}.
\]
Its contribution is
\[
 \ll\frac{\sqrt K(\log K)^{3/2}}T
 =\frac{(\log K)^{B_{\rm tail}+3/2}}{\sqrt K},
\]
and hence the argument has proved the more precise intermediate estimate
\begin{equation}\label{eq:ram-tail-two-term}
 \left\|\sum_{q>Q}\frac{c_q(n)}{q^2}\right\|_{\ell^2(n\le2K)}
 \ll \frac{\sqrt K}{Q}
      +E_K,\qquad
 E_K:=\frac{(\log K)^{B_{\rm tail}+3/2}}{\sqrt K},
 \quad 2\le Q\le T.
\end{equation}
The same estimate holds for finite tails, and the same divisor calculation
also gives, uniformly in the upper endpoint,
\begin{equation}\label{eq:ram-tail-terminal}
 \left\|\sum_{x<q\le R}\frac{c_q(n)}{q^2}\right\|_{\ell^2(n\le2K)}
 \ll\frac{\sqrt K(\log K)^{3/2}}{x},
 \qquad x\ge T.
\end{equation}
For all sufficiently large
\(K\), \(E_K=o(1)\), whereas
\(\sqrt K/Q\ge1\) in the stated range; the bounded remaining values are
absorbed into the constant.  Letting \(R\to\infty\) proves the infinite
tail.  Notice that the larger range
\(Q\le K/(\log K)^{B_{\rm tail}}\) would not follow from this argument: at its
upper endpoint the displayed terminal error is too large by a factor
\((\log K)^{3/2}\).  The range stated here is exactly the one used below.
\end{proof}

\subsection{A smooth near-resonant estimate}

Fix \(0<\varepsilon<1/2\), \(A\ge1\), and put \(a=A/L\).
Fix Gevrey-order-two functions \(\eta,\beta\) on \([0,\infty)\), with
\[
 \eta=0\ \text{on }[0,1],\quad \eta=1\ \text{on }[2,\infty),
 \qquad
 \beta=1\ \text{on }[0,1/2],\quad \beta=0\ \text{on }[1,\infty),
\]
and set
\[
 \rho(v)=\rho_{a,\varepsilon}(v)
 :=\eta(|v|/a)\beta(|v|/\varepsilon).
\]
Thus \(\rho\) is even, vanishes when
\(|v|\le a\) or \(|v|\ge\varepsilon\), equals one when
\[
 2a\le|v|\le\varepsilon/2,
\]
and satisfies
\begin{equation}\label{eq:gevrey}
 \norm{\rho^{(j)}}_\infty
 \le C_\varepsilon^{j+1}(j!)^2a^{-j}.
\end{equation}
Define
\begin{equation}\label{eq:Ynear}
 Y^{\mathrm{near}}_{N,A,\varepsilon}(\alpha)=
 \sum_{\substack{p\le N\\(p,q_p)=1}}
 \frac{V(N\delta_p)}{p\delta_p}\rho(p\delta_p).
\end{equation}

\begin{lemma}[Uniform inner-cutoff multipliers]
\label{lem:near-multiplier}
Suppose that \(0<a\le\varepsilon/4\), and put
\[
 W_a(v)=
 \begin{cases}
  \rho_{a,\varepsilon}(v)/v,&v\ne0,\\
  0,&v=0,
 \end{cases}
 \qquad
 J_a(t)=\int_{\R}W_a(v)e(-tv)\dd v.
\]
Then \(W_a\in C_c^\infty(\R)\).  For every integer \(m\ge1\) and
every \(t\ne0\),
\begin{equation}\label{eq:J-uniform-bounds}
 |J_a(t)|+|tJ_a'(t)|
 \le
 C_\varepsilon^{m+1}(m!)^2
 \min\left\{|t|,1,(a|t|)^{-m}\right\}.
\end{equation}
In particular, \(\sup_{0<a\le\varepsilon/4,t\in\R}|J_a(t)|
\ll_\varepsilon1\).  Moreover,
\begin{equation}\label{eq:J-uniform-envelope}
 \int_0^\infty
 \sup_{1\le r\le2}r|J_a'(rt)|\dd t
 \ll_\varepsilon1,
\end{equation}
uniformly for \(0<a\le\varepsilon/4\), and, for every integer \(j\ge0\),
\begin{equation}\label{eq:W-derivatives}
 \norm{W_a^{(j)}}_{L^2(\R)}
 \le C_\varepsilon^{j+1}(j!)^2a^{-j-1/2}.
\end{equation}
All constants displayed here are independent of \(a\).
\end{lemma}

\begin{proof}
Set
\[
 z_\varepsilon(v)=\beta(|v|/\varepsilon),\qquad
 h(u)=1-\eta(|u|),\qquad h_a(v)=h(v/a).
\]
Because \(a\le\varepsilon/4\), the support of \(h_a\) is contained in
\(\{|v|\le2a\}\), where \(z_\varepsilon=1\).  On that support,
\(z_\varepsilon-h_a=\eta(|v|/a)=\rho_{a,\varepsilon}\); off that
support, \(h_a=0\) and \(\eta(|v|/a)=1\).  Hence
\begin{equation}\label{eq:rho-two-profiles}
 \rho_{a,\varepsilon}=z_\varepsilon-h_a.
\end{equation}
Both terms on the right are fixed smooth compactly supported profiles,
the second followed by dilation by \(a\).

For every fixed integer \(M\ge2\), integration by parts \(M\) times
and scaling give, for \(t\ge0\),
\begin{equation}\label{eq:rho-Fourier-envelope}
 \sup_{1\le r\le2}|\widehat\rho_{a,\varepsilon}(rt)|
 \le C_{\varepsilon,M}\left\{
 \varepsilon(1+\varepsilon t)^{-M}
 +a(1+at)^{-M}\right\}.
\end{equation}
Since
\[
 J_a'(t)=-2\pi\ii\,\widehat\rho_{a,\varepsilon}(t),
\]
the integral of the right side of
\eqref{eq:rho-Fourier-envelope} over \(t>0\) is
\(O_{\varepsilon,M}(1)\).  This proves
\eqref{eq:J-uniform-envelope} with the requested explicit envelope.

We next prove \eqref{eq:J-uniform-bounds}.  The function
\(\rho_{a,\varepsilon}\) is even, so \(W_a\) is odd and
\[
 J_a(t)=-2\ii\int_0^\infty
 \rho_{a,\varepsilon}(v)\frac{\sin(2\pi tv)}v\dd v.
\]
Since
\[
 \frac{\sin(2\pi tv)}v\dd v=\dd\Si(2\pi tv),
\]
Stieltjes integration by parts yields
\[
 |J_a(t)|
 \le2\sup_{x\in\R}|\Si(x)|
 \TV_{(0,\infty)}(\rho_{a,\varepsilon})
 \ll_\varepsilon1.
\]
The total variation is uniform in \(a\), either directly from the two
transition regions or from \eqref{eq:rho-two-profiles}.  Similarly,
Fourier integration by parts gives
\[
 |tJ_a'(t)|
 =2\pi|t\widehat\rho_{a,\varepsilon}(t)|
 \le\norm{\rho_{a,\varepsilon}'}_1
 \ll_\varepsilon1.
\]
At the origin, \(J_a(0)=0\) by oddness and
\[
 |J_a'(t)|\le2\pi\norm{\rho_{a,\varepsilon}}_1
 \ll_\varepsilon1.
\]
Thus
\[
 |J_a(t)|+|tJ_a'(t)|\ll_\varepsilon |t|
 \qquad (|t|\le1).
\]

It remains to record the high-frequency estimates with their dependence
on the order.  Leibniz's rule, the Gevrey bounds for the two transition
profiles, and their supports imply, for \(m\ge1\),
\begin{equation}\label{eq:W-L1-derivatives}
 \norm{W_a^{(m)}}_1
 \le C_\varepsilon^{m+1}(m!)^2a^{-m},
 \qquad
 \norm{\rho_{a,\varepsilon}^{(m+1)}}_1
 \le C_\varepsilon^{m+1}(m!)^2a^{-m}.
\end{equation}
For the first estimate, expand
\[
 W_a^{(m)}
 =\sum_{k=0}^m
 \binom mk\rho_{a,\varepsilon}^{(k)}
 \left(\frac1v\right)^{(m-k)}.
\]
For \(k=0\), integration over \(a<|v|<\varepsilon\) gives
\(O(m!a^{-m})\).  For \(k\ge1\), the derivative of \(\rho\) is
supported in
\[
 a<|v|<2a,\qquad \varepsilon/2<|v|<\varepsilon.
\]
The inner region contributes at most
\[
 C_\varepsilon^{m+1}(m!)^2
 a^{-k}a^{-(m-k+1)}a
 \ll C_\varepsilon^{m+1}(m!)^2a^{-m},
\]
and the outer region is smaller.  The second estimate in
\eqref{eq:W-L1-derivatives} follows in the same way by scaling.
Integrating by parts \(m\) times in \(J_a\), and \(m+1\) times in
\(\widehat\rho\), now gives
\[
 |J_a(t)|+|tJ_a'(t)|
 \le C_\varepsilon^{m+1}(m!)^2(a|t|)^{-m}.
\]
Together with the small- and intermediate-frequency bounds, this proves
\eqref{eq:J-uniform-bounds}.

The same Leibniz expansion in \(L^2\) gives
\[
 \left\|
 \rho_{a,\varepsilon}^{(k)}
 \left(\frac1v\right)^{(j-k)}
 \right\|_2
 \le C_\varepsilon^{j+1}(j!)^2a^{-j-1/2}.
\]
For \(k=0\), integrate \(|v|^{-2(j+1)}\) from \(a\) to
\(\varepsilon\).  For \(k\ge1\), the inner transition region contributes
\[
 a^{-k}a^{-(j-k+1)}a^{1/2}=a^{-j-1/2},
\]
and the outer region is smaller.  Summing over \(k\) proves
\eqref{eq:W-derivatives}.
\end{proof}

\begin{proposition}[Near-resonant square function]\label{prop:near}
For fixed \(A\ge1\) and \(0<\varepsilon<1/2\),
\begin{equation}\label{eq:near}
 \limsup_{N\to\infty}
 \frac{\norm{Y^{\mathrm{near}}_{N,A,\varepsilon}}_2^2}{L^2}
 \ll_\varepsilon\frac1A.
\end{equation}
\end{proposition}

\begin{proof}
For fixed \(A,\varepsilon\), the relation \(a=A/L\le\varepsilon/4\)
holds for all sufficiently large \(N\), which is the range relevant to
the limsup.  Put \(W=W_a\), \(J=J_a\), and
\[
 b_p(\alpha)=
 \1_{\{(p,q_p)=1\}}\frac{\rho(p\delta_p)}{p\delta_p}.
\]
For each \(p\), there is the exact locally finite periodization
\begin{equation}\label{eq:bp-cell-periodization}
 b_p(\alpha)=
 \sum_{\substack{q\in\Z\\(p,q)=1}}
 W\left(p^2\left(\alpha-\frac qp\right)\right).
\end{equation}
Indeed, the support of the summand centered at \(q/p\) has radius
\(\varepsilon/p^2\).  Distinct centers are separated by \(1/p\), and
\(2\varepsilon/p^2<1/p\).  Thus the supports are disjoint.  On such a
support, \(|p\alpha-q|<\varepsilon/p<1/2\), so \(q=q_p\), and
\(p^2(\alpha-q/p)=p\delta_p\).  Since
\(W\in C_c^\infty(\R)\) and vanishes with all derivatives at its cutoff
boundaries, \eqref{eq:bp-cell-periodization} also proves that \(b_p\)
is a periodic \(C^\infty\) function.  Restricting to reduced \(q\)'s
removes entire disjoint smooth cells and creates no jump.

Changing variables in those cells gives, for every \(n\in\Z\),
\begin{equation}\label{eq:near-bhat}
 \begin{aligned}
 \wh b_p(n)
 &=\frac1{p^2}
 \sum_{\substack{q\bmod p\\(p,q)=1}}e(-nq/p)
 \int_{\R}W(v)e(-nv/p^2)\dd v\\
 &=\frac{c_p(-n)}{p^2}J(n/p^2).
 \end{aligned}
\end{equation}
In particular \(J(0)=0\), since \(W\) is odd.

We begin with the constant coefficient \(v_0=1/12\) in
\eqref{eq:Vfourier}.  Here and below in this proof, \(n\asymp K\)
means \(K<n\le2K\).  Put
\[
 Q_K=\sqrt{aK},\qquad R_K=\sqrt{\varepsilon K}.
\]
Lemma~\ref{lem:near-multiplier} and
\[
 p\frac{\dd}{\dd p}J(n/p^2)
 =-2\frac n{p^2}J'(n/p^2)
\]
give, for every fixed integer \(m\ge1\),
\begin{equation}\label{eq:Jbounds}
 |J(n/p^2)|+
 p\left|\frac{\dd}{\dd p}J(n/p^2)\right|
 \ll_{m,\varepsilon}
 \begin{cases}
 (p/Q_K)^{2m},&p<Q_K,\\
 1,&Q_K\le p\le R_K,\\
 K/p^2,&p>R_K.
 \end{cases}
\end{equation}
For the last line, use the small-\(t\) part of
\eqref{eq:J-uniform-bounds} when \(n/p^2\le1\).  If
\(1<n/p^2\le2/\varepsilon\), the uniform bound is at most
\(C_\varepsilon K/p^2\): indeed \(p^2<n\le2K\), so
\(K/p^2\ge1/2\), and the uniform \(O_\varepsilon(1)\) bound is
absorbed by the displayed quantity.
We now write every boundary term in the multiplier summation.  Let
\[
 \mathcal H_K=\ell^2\{n\in\Z:K<n\le2K\},\qquad
 u_p=\left(\frac{c_p(-n)}{p^2}\right)_{n\asymp K},
 \qquad
 w_x=\bigl(J(n/x^2)\bigr)_{n\asymp K}.
\]
For \(2\le U<\infty\), define the finite tail
\[
 \mathcal R_{x,U}=\sum_{x<p\le U}u_p.
\]
Writing \(\odot\) for coordinatewise multiplication, summation by parts
in each coordinate gives the exact finite-dimensional identity
\begin{equation}\label{eq:near-vector-Abel}
 \sum_{2<p\le U}u_p\odot w_p
 =
 \mathcal R_{2,U}\odot w_2+
 \int_2^U\mathcal R_{x,U}\odot
 \frac{\dd w_x}{\dd x}\dd x.
\end{equation}
Indeed, \(x\mapsto\mathcal R_{x,U}\) has jump \(-u_p\) at \(p\)
and \(\mathcal R_{U,U}=0\), so Stieltjes integration by parts is
\[
 -\int_{(2,U]}w_x\odot\dd\mathcal R_{x,U}
 =\mathcal R_{2,U}\odot w_2+
  \int_2^U\mathcal R_{x,U}\odot\dd w_x.
\]
This proves \eqref{eq:near-vector-Abel}, including both endpoint terms.

Let
\[
 T=\frac K{(\log K)^{B_{\rm tail}}},\qquad
 E_K=\frac{(\log K)^{B_{\rm tail}+3/2}}{\sqrt K}.
\]
Since Ramanujan sums are real and even, the finite-tail form of
\eqref{eq:ram-tail-two-term} gives, uniformly for
\(2\le x\le U\le T\),
\begin{equation}\label{eq:near-finite-tail}
 \norm{\mathcal R_{x,U}}_{\mathcal H_K}
 \ll\frac{\sqrt K}{x}+E_K.
\end{equation}
For a fixed \(m\ge1\), put
\[
 \omega_{K,m}(x)=
 \begin{cases}
  (x/Q_K)^{2m},&x<Q_K,\\
  1,&x\ge Q_K.
 \end{cases}
\]
Equation \eqref{eq:Jbounds} yields
\begin{equation}\label{eq:near-weight-envelope}
 \norm{w_x}_{\ell^\infty(n\asymp K)}
 +x\norm{\frac{\dd w_x}{\dd x}}_{\ell^\infty(n\asymp K)}
 \ll_{\varepsilon,m}\omega_{K,m}(x).
\end{equation}
The uniform envelope in Lemma~\ref{lem:near-multiplier} also gives
\begin{equation}\label{eq:J-total-variation}
 \int_0^\infty
 \norm{\frac{\dd w_x}{\dd x}}_{\ell^\infty(n\asymp K)}
 \dd x\ll_\varepsilon1.
\end{equation}
To check the change of variables, write \(n=rK\), \(1\le r\le2\),
and \(t=K/x^2\); the left side is at most a constant times
\[
 \int_0^\infty\sup_{1\le r\le2}r|J'(rt)|\dd t,
\]
which is \eqref{eq:J-uniform-envelope}.

Assume first \(K\ge16/a\), so \(Q_K\ge4\), and set
\(U_0=\min\{N,T\}\).  Equations
\eqref{eq:near-vector-Abel}--\eqref{eq:J-total-variation} give
\begin{align}
 \left\|\sum_{2<p\le U_0}u_p\odot w_p\right\|_{\mathcal H_K}
 &\le
 \left(\frac{\sqrt K}{2}+E_K\right)\norm{w_2}_\infty
 +\int_2^{U_0}
 \left(\frac{\sqrt K}{x}+E_K\right)
 \norm{\frac{\dd w_x}{\dd x}}_\infty\dd x \notag\\
 &\ll_{\varepsilon,m}
 \sqrt K\left\{
 \frac{\omega_{K,m}(2)}2+
 \int_2^\infty\frac{\omega_{K,m}(x)}{x^2}\dd x
 \right\}+E_K \notag\\
 &\ll_{\varepsilon,m}\frac{\sqrt K}{Q_K}+E_K.
 \label{eq:near-Abel-low}
\end{align}
The last line follows from
\[
 Q_K^{-2m}\int_2^{Q_K}x^{2m-2}\dd x
 +\int_{Q_K}^\infty x^{-2}\dd x\ll_m Q_K^{-1},
\qquad
\omega_{K,m}(2)\ll_m Q_K^{-1}.
\]
If \(T<N\), apply \eqref{eq:near-vector-Abel} on \((T,N]\).
The terminal tail estimate \eqref{eq:ram-tail-terminal} and
\eqref{eq:near-weight-envelope} give
\begin{align}
 \left\|\sum_{T<p\le N}u_p\odot w_p\right\|_{\mathcal H_K}
 &\ll_\varepsilon
 \sqrt K(\log K)^{3/2}
 \left(\frac1T+\int_T^\infty\frac{\dd x}{x^2}\right)\notag\\
 &\ll_\varepsilon
 \frac{\sqrt K(\log K)^{3/2}}T=E_K.
 \label{eq:near-Abel-terminal}
\end{align}
The finitely many terms \(p\le2\) satisfy the same
\(O_{\varepsilon,m}(\sqrt K/Q_K)\) bound by the first line of
\eqref{eq:Jbounds}.
Consequently, whenever \(K\ge16/a\), uniformly over the dyadic block,
\begin{equation}\label{eq:near-constant-block}
 \left\|
 \sum_{p\le N}\frac{c_p(n)}{p^2}J(n/p^2)
 \right\|_{\ell^2(n\asymp K)}
 \ll_\varepsilon \frac{\sqrt K}{Q_K}+E_K
 \ll_\varepsilon a^{-1/2}.
\end{equation}
The assertion about \(E_K\) is uniform, not merely pointwise.  Indeed,
for all sufficiently large \(N\), the function
\(x\mapsto(\log x)^{B_{\rm tail}+3/2}x^{-1/2}\) is decreasing throughout
\(x\ge16/a\), and therefore
\begin{equation}\label{eq:near-E-uniform}
 \sup_{K\ge16/a}E_K
 \le C\sqrt a\,\bigl(\log(16/a)\bigr)^{B_{\rm tail}+3/2}=o(1).
\end{equation}

We spell out both ends of the frequency sum.  For
\(1\le n\le16/a\), absolute convergence, the uniform bound for \(J\),
and the divisor formula give
\begin{equation}\label{eq:near-initial-pointwise}
 \left|\sum_{p\le N}\frac{c_p(-n)}{p^2}J(n/p^2)\right|
 \ll\sum_{p\ge1}\frac{|c_p(n)|}{p^2}
 \ll\sigma_{-1}(n).
\end{equation}
Consequently, using
\(\sigma_{-1}(n)\ll1+\log\log(3n)\), these frequencies contribute at
most
\begin{equation}\label{eq:near-initial-sum}
 O_\varepsilon\left(
 a^{-1}(1+\log\log(1/a))^2\right)=o(L/a)
\end{equation}
to the squared norm.

At the other endpoint suppose that \(Q_K>N\), equivalently
\(K>N^2/a\).  For all sufficiently large \(N\), this implies \(N<T\),
because \(x/(\log x)^{B_{\rm tail}}\) is then increasing for
\(x\ge N^2/a\), and
\[
 T\ge\frac{N^2/a}{(\log(N^2/a))^{B_{\rm tail}}}>N.
\]
Apply \eqref{eq:near-vector-Abel} with \(U=N\).  Every \(x\le N\)
then lies below \(Q_K\), and hence
\begin{equation}\label{eq:near-low-envelope-tail}
 \norm{w_x}_\infty+x\norm{w_x'}_\infty
 \ll_{\varepsilon,m}(x/Q_K)^{2m}.
\end{equation}
Using \eqref{eq:near-finite-tail} in the exact Abel identity gives
\begin{align}
 \left\|\sum_{2<p\le N}u_p\odot w_p\right\|_{\mathcal H_K}
 &\ll_{\varepsilon,m}
 \sqrt K Q_K^{-2m}
 \left(1+\int_2^N x^{2m-2}\dd x\right)
 +E_K\left(\frac N{Q_K}\right)^{2m}\notag\\
 &\ll_{\varepsilon,m}
 \frac{\sqrt K}{Q_K}
 \left(\frac N{Q_K}\right)^{2m-1}
 +E_K\left(\frac N{Q_K}\right)^{2m}.
 \label{eq:near-terminal-block-detail}
\end{align}
The terms \(p\le2\) obey the same first bound directly.  By
\eqref{eq:near-E-uniform}, \(E_K\le a^{-1/2}=\sqrt K/Q_K\) for all
large \(N\).  Since \(N/Q_K<1\), we conclude that
\begin{equation}\label{eq:near-terminal-block}
 \left\|\sum_{p\le N}\frac{c_p(-n)}{p^2}J(n/p^2)
 \right\|_{\ell^2(n\asymp K)}
 \ll_{\varepsilon,m}
 a^{-1/2}\left(\frac N{Q_K}\right)^{2m-1}.
\end{equation}
On the dyadic blocks \(K_r=2^rN^2/a\), the square of the last bound is
at most
\[
 \frac1a\left(\frac{N^2}{aK_r}\right)^{2m-1}
 =\frac1a\cdot2^{-r(2m-1)}.
\]
Thus the complete squared tail \(K>N^2/a\) is
\(O_{\varepsilon,m}(a^{-1})\).  Between \(16/a\) and \(N^2/a\)
there are \(O(L)\) dyadic blocks, and each contributes
\(O_\varepsilon(a^{-1})\) by \eqref{eq:near-constant-block}.
Together with \eqref{eq:near-initial-sum}, this proves that the constant
mode has squared norm
\begin{equation}\label{eq:near-constant}
 \ll_\varepsilon L/a=\frac{L^2}{A}.
\end{equation}

Negative frequencies contribute the same amount by conjugation, and the
zero frequency vanishes because \(J(0)=0\).

Now fix \(\ell\ne0\).  Since \(N,\ell,q_p\) are integers,
\[
 e(\ell N\delta_p)=e(\ell Np\alpha).
\]
Differentiating the disjoint-cell representation
\eqref{eq:bp-cell-periodization} and changing variables in every cell
gives, for every integer \(j\ge0\),
\begin{align}
 \norm{b_p^{(j)}}_2^2
 &=\varphi(p)p^{4j-2}\norm{W^{(j)}}_2^2,\notag\\
 \norm{b_p^{(j)}}_2
 &\le C_\varepsilon^{j+1}(j!)^2
 p^{2j-1/2}a^{-j-1/2}\notag\\
 &=C_\varepsilon^{j+1}(j!)^2
 \left(\frac{p^2}{a}\right)^j(ap)^{-1/2}.
 \label{eq:bp-derivative}
\end{align}
In particular,
\begin{equation}\label{eq:bp-L2}
 \norm{b_p}_2^2\ll_\varepsilon\frac1{ap}.
\end{equation}
Whenever \(b_p(\alpha)\ne0\),
\[
 \left|\alpha-\frac{q_p}p\right|
 =\frac{|p\delta_p|}{p^2}<\frac1{2p^2},
\]
and \((p,q_p)=1\).  Legendre's criterion therefore makes \(p\) a
continued-fraction denominator of \(\alpha\).  Since such denominators
satisfy \(Q_{n+2}\ge2Q_n\), only an absolute bounded number of active
\(p\)'s can lie in a dyadic interval.  Let
\[
 P_0=N\exp(-L^{1/3}).
\]
The same recurrence shows that the total number of active denominators
in \((P_0,N]\) is \(O(L^{1/3})\).  Pointwise Cauchy--Schwarz and
\eqref{eq:bp-L2} therefore give
\begin{align}
 \left\|\sum_{P_0<p\le N}b_p(\alpha)e(\ell Np\alpha)\right\|_2^2
 &\ll L^{1/3}\sum_{P_0<p\le N}\norm{b_p}_2^2\notag\\
 &\ll\frac{L^{1/3}}a\sum_{P_0<p\le N}\frac1p
 \ll\frac{L^{2/3}}a=o(L/a).
 \label{eq:near-highp}
\end{align}

For \(p\le P_0\), divide into dyadic blocks
\(\mathcal P_s=(P_s/2,P_s]\), and set
\[
 F_{s,\ell}(\alpha)=
 \sum_{p\in\mathcal P_s}b_p(\alpha)e(\ell Np\alpha).
\]
The same pointwise sparsity and \eqref{eq:bp-L2} give
\begin{equation}\label{eq:near-Qblock}
 \norm{F_{s,\ell}}_2^2
 \ll\sum_{p\in\mathcal P_s}\norm{b_p}_2^2
 \ll_\varepsilon\frac1a.
\end{equation}
Let \(\Pi_{s,\ell}\) be Fourier projection onto the signed annulus
\[
 \mathcal A_{s,\ell}=
 \left\{n\in\Z:\
 \sgn(n)=\sgn(\ell),\quad
 \frac{|\ell|NP_s}{4}\le|n|\le2|\ell|NP_s\right\}.
\]
We use the following exact leakage inequality.  If \(f\in H^j(\T)\),
\(m_0\in\Z\), and every integer outside a frequency set
\(\mathcal A\) has distance at least \(D\) from \(m_0\), then
Plancherel gives
\begin{align}
 \norm{(I-\Pi_{\mathcal A})(e(m_0\alpha)f)}_2^2
 &=\sum_{\substack{k\in\Z\\m_0+k\notin\mathcal A}}
 |\widehat f(k)|^2\notag\\
 &\le D^{-2j}\sum_{k\in\Z}|k|^{2j}|\widehat f(k)|^2
 =(2\pi D)^{-2j}\norm{f^{(j)}}_2^2.
 \label{eq:near-Fourier-leakage}
\end{align}
For \(p\in(P_s/2,P_s]\), the carrier \(m_0=\ell Np\) lies in
\(\mathcal A_{s,\ell}\), and its distance from the complement is at
least \(D_{s,\ell}=|\ell|NP_s/4\).  Applying
\eqref{eq:near-Fourier-leakage}, the triangle inequality, and
\eqref{eq:bp-derivative}, we obtain
\begin{align}
 \norm{(I-\Pi_{s,\ell})F_{s,\ell}}_2
 &\le(2\pi D_{s,\ell})^{-j}
 \sum_{p\in\mathcal P_s}\norm{b_p^{(j)}}_2\notag\\
 &\le C_\varepsilon^{j+1}(j!)^2|\ell|^{-j}
 \left(\frac{P_s}{Na}\right)^j
 \left(\frac{P_s}{a}\right)^{1/2}.
 \label{eq:near-leakage}
\end{align}
There is no unrecorded dependence on \(j\) in this estimate.  Choose
\(j=\ceil{2L^{2/3}}\).  Since
\[
 \log\frac{P_s}{Na}\le-L^{1/3}+\log(L/A),
\]
\(|\ell|^{-j}\le1\), and \(\log(j!)\le j\log j\), the logarithm of
the last line of \eqref{eq:near-leakage} is at most
\[
 -jL^{1/3}+j\log(L/A)+2j\log j+O_\varepsilon(j)
 +\frac12\log(P_s/a)
 =-\frac32L+O_\varepsilon(L^{2/3}\log L).
\]
It follows that
\begin{equation}\label{eq:near-leakage-small}
 \sup_{s,\,\ell\ne0}
 \norm{(I-\Pi_{s,\ell})F_{s,\ell}}_2=o(N^{-1}).
\end{equation}

The annuli \(\mathcal A_{s,\ell}\), with \(P_s\) dyadic, have overlap
multiplicity at most an absolute constant \(B\).  Consequently,
frequency by frequency and then by Plancherel,
\begin{align}
 \left\|\sum_s\Pi_{s,\ell}F_{s,\ell}\right\|_2^2
 &=\sum_{n\in\Z}
 \left|\sum_{s:n\in\mathcal A_{s,\ell}}
 \widehat F_{s,\ell}(n)\right|^2\notag\\
 &\le B\sum_s\sum_{n\in\mathcal A_{s,\ell}}
 |\widehat F_{s,\ell}(n)|^2
 \le B\sum_s\norm{F_{s,\ell}}_2^2
 \ll_\varepsilon L/a.
 \label{eq:near-carrier-square}
\end{align}
There are \(O(L)\) blocks, so \eqref{eq:near-leakage-small} also gives
\[
 \sum_s\norm{(I-\Pi_{s,\ell})F_{s,\ell}}_2=o(1).
\]
The triangle inequality and \eqref{eq:near-carrier-square} therefore
give
\begin{equation}\label{eq:near-lowp-shifted}
 \left\|\sum_{p\le P_0}b_p(\alpha)e(\ell Np\alpha)\right\|_2^2
 \ll_\varepsilon L/a,
\end{equation}
uniformly for \(\ell\ne0\).
Together with \eqref{eq:near-highp}, this is uniform in
\(\ell\ne0\).  Finally, the absolutely convergent Fourier series
\(V(x)=\sum_{\ell\in\Z}v_\ell e(\ell x)\) gives
\[
 Y^{\mathrm{near}}_{N,A,\varepsilon}
 =\sum_{\ell\in\Z}v_\ell
 \sum_{p\le N}b_p(\alpha)e(\ell N\delta_p).
\]
Use \eqref{eq:near-constant} for \(\ell=0\), the shifted estimate for
\(\ell\ne0\), Minkowski's inequality, and
\(\sum_{\ell\ne0}|v_\ell|<\infty\).  We obtain
\[
 \norm{Y^{\mathrm{near}}_{N,A,\varepsilon}}_2^2
 \ll_\varepsilon\frac La=\frac{L^2}{A},
\]
which proves \eqref{eq:near}.
\end{proof}

\subsection{The fixed-away estimate}

\begin{proposition}[Fixed-away shots]\label{prop:far}
For every fixed \(0<\varepsilon<1/2\),
\begin{equation}\label{eq:far}
 \left\|
 \sum_{\substack{p\le N,\ (p,q_p)=1\\
                  |p\delta_p|\ge\varepsilon}}
 \frac{V(N\delta_p)}{p\delta_p}
 \right\|_2=o(\log N).
\end{equation}
More precisely, for every \(0<\varepsilon_0<\varepsilon_1<1/2\),
\begin{equation}\label{eq:far-uniform}
 \lim_{N\to\infty}
 \sup_{t\in[\varepsilon_0,\varepsilon_1]}
 \frac1{\log N}
 \left\|
 \sum_{\substack{p\le N,\ (p,q_p)=1\\|p\delta_p|\ge t}}
 \frac{V(N\delta_p)}{p\delta_p}
 \right\|_2=0.
\end{equation}
\end{proposition}

\begin{proof}
Fix temporarily \(0<\delta<\varepsilon/2\), and take an even smooth
function \(\chi=\chi_{\varepsilon,\delta}\) such that
\[
 \chi(v)=0\quad(|v|\le\varepsilon-\delta),\qquad
 \chi(v)=1\quad(|v|\ge\varepsilon).
\]
Let \(Y_N^\chi\) be \eqref{eq:defY} with the extra factor
\(\chi(p\delta_p)\), and define its full periodization by
\[
 Z_N^\chi(\alpha)=
 \sum_{p\le N}\frac1p\,
 \PV\sum_{\substack{q\in\Z\\(p,q)=1}}
 \frac{V(N(p\alpha-q))}{p\alpha-q}\,
 \chi\bigl(p(p\alpha-q)\bigr).
\]
Because \(1-\chi\) is supported in \(|p(p\alpha-q)|<\varepsilon<1/2\),
there is at most one \(q\) in its support, necessarily \(q_p\).
Therefore
\begin{equation}\label{eq:far-window}
 Y_N^\chi-Z_N^\chi=Y_N-Z_{N,N},
\end{equation}
and Proposition~\ref{prop:window} gives
\begin{equation}\label{eq:far-window-L2}
 \norm{Y_N^\chi-Z_N^\chi}_2^2\ll L\log L=o(L^2).
\end{equation}

Set
\begin{equation}\label{eq:Rchi}
 R_\chi(t)=\PV\int_{\R}\frac{\chi(v)}v e(-tv)\dd v.
\end{equation}
Put \(\kappa_\chi=1-\chi\).  This function is even and compactly
supported.  For \(t\ne0\), the Fourier transform of
\(\PV(1/v)\) gives
\begin{equation}\label{eq:Rchi-decomposition}
 R_\chi(t)=-\ii\pi\sgn(t)-
 \PV\int_{\R}\frac{\kappa_\chi(v)}v e(-tv)\dd v.
\end{equation}
The second term in \eqref{eq:Rchi-decomposition} tends to zero as
\(t\to0\).  Indeed, after pairing \(v\) and \(-v\), its absolute
value is at most
\[
 2\int_0^C\kappa_\chi(v)
 \frac{|\sin(2\pi tv)|}{v}\dd v=O_\chi(|t|).
\]
Thus the symmetric value and the two one-sided limits are
\begin{equation}\label{eq:Rchi-jump}
 R_\chi(0)=0,\qquad
 R_\chi(0^+)=-\ii\pi,qquad R_\chi(0^-)=\ii\pi.
\end{equation}
For \(t\ne0\), differentiation in
\eqref{eq:Rchi-decomposition} gives
\[
 R_\chi'(t)=2\pi\ii\widehat{\kappa_\chi}(t).
\]
Dirichlet's integral applied to \eqref{eq:Rchi} shows that
\(R_\chi(t)\to0\) at both infinities.  Since
\(\widehat{\kappa_\chi}\) is Schwartz, integration from the
appropriate infinity shows that \(R_\chi\) is smooth and rapidly
decreasing on each open half-line.  More precisely, for all
\(r,J\ge0\),
\begin{equation}\label{eq:Rchi-Schwartz}
 |R_\chi^{(r)}(t)|
 \le C_{r,J,\chi}(1+|t|)^{-J}\qquad(t\ne0),
\end{equation}
where for \(r=0\) the two half-lines are treated separately.  In
particular, \(R_\chi\), with the value in \eqref{eq:Rchi-jump}, is a
regulated BV function with one jump; it is not being regarded as a
Schwartz function on all of \(\R\).

\medskip
\noindent\textbf{Auxiliary scaled-multiplier estimate.}
Let \(R=R_\chi\).  If \(p,p'\in(Q,2Q]\) and \(a,a'\in\R\), then,
for every \(J\ge0\),
\begin{equation}\label{eq:scaled-BV-general}
 \begin{split}
 W(t)&=R\left(\frac{t-a}{p^2}\right)
       \overline{R\left(\frac{t-a'}{p'^2}\right)},\\
 \norm{W}_\infty+\TV_{\R}(W)
 &\ll_{J,\chi}
 \left(1+\frac{|a-a'|}{Q^2}\right)^{-J}.
 \end{split}
\end{equation}
Here the values at \(a\) and \(a'\) are induced by \(R(0)=0\).
Moreover, suppose that \(\ell\ne0\),
\(a=\ell Np\), \(a'=\ell Np'\), and put
\[
 H=\frac{|\ell|N}{Q},\qquad
 I_{\ell,Q}=\left\{t:\ \sgn(t)=\sgn(\ell),\quad
 \frac{|\ell|NQ}{2}\le|t|\le3|\ell|NQ\right\},
 \qquad E_{\ell,Q}=\R\setminus I_{\ell,Q}.
\]
Then
\begin{equation}\label{eq:scaled-BV-projected}
 \norm{\1_{E_{\ell,Q}}W}_\infty+
 \TV_{\R}(\1_{E_{\ell,Q}}W)
 \ll_{J,\chi}H^{-J}.
\end{equation}
The convention at the two endpoints of \(I_{\ell,Q}\) is immaterial.

\smallskip
\noindent\emph{Proof of the auxiliary estimate.}
Write \(E_M(x)=(1+|x|)^{-M}\).  Uniformly for
\(1/4\le\lambda\le4\), \(d\in\R\), and \(M>J+1\),
\begin{equation}\label{eq:two-envelope-convolution}
 \sup_x E_M(x)E_M(\lambda x-d)+
 \int_{\R}E_M(x)E_M(\lambda x-d)\dd x
 \ll_{J,M}(1+|d|)^{-J}.
\end{equation}
To verify this, split the line into \(|x|\ge|d|/8\) and its
complement.  On the first set one extracts
\((1+|d|)^{-J}\) from the first factor and integrates the remaining
envelope.  On the complement,
\(|\lambda x-d|\ge|d|/2\), and one instead extracts the same power
from the second factor.  The supremum estimate is identical.

Away from \(a,a'\), ordinary differentiation gives
\[
 W'(t)=\frac1{p^2}R'\left(\frac{t-a}{p^2}\right)
 \overline{R\left(\frac{t-a'}{p'^2}\right)}
 +\frac1{p'^2}R\left(\frac{t-a}{p^2}\right)
 \overline{R'\left(\frac{t-a'}{p'^2}\right)}.
\]
In the first integral use \(t=a+p^2x\), and in the second use
\(t=a'+p'^2x\).  Since \(p^2/p'^2\in[1/4,4]\),
\eqref{eq:Rchi-Schwartz} and
\eqref{eq:two-envelope-convolution} bound both
\(\int_{\R\setminus\{a,a'\}}|W'(t)|\dd t\) and
\(\norm{W}_\infty\) by the right side of
\eqref{eq:scaled-BV-general}.

It remains to count the point jumps.  If \(a\ne a'\), then the total
variation contributed at \(a\), including the assigned value
\(W(a)=0\), is
\[
 2\pi\left|R\left(\frac{a-a'}{p'^2}\right)\right|;
\]
the analogous contribution at \(a'\) is
\(2\pi|R((a'-a)/p^2)|\).  Both have the required separation decay.
If \(a=a'\), the two one-sided products are bounded and the assigned
value is zero, so the total point contribution is at most \(2\pi^2\).
This proves \eqref{eq:scaled-BV-general}, including unequal scales and
all possible coincidences of the jump points.

For \eqref{eq:scaled-BV-projected}, both centers lie strictly inside
\(I_{\ell,Q}\), while for every \(t\in E_{\ell,Q}\),
\[
 1+\frac{|t-a|}{p^2}\gg1+H,
 \qquad
 1+\frac{|t-a'|}{p'^2}\gg1+H.
\]
Thus multiplication by \(\1_{E_{\ell,Q}}\) kills the original jump
points.  On \(E_{\ell,Q}\), use the displayed derivative formula,
take the supremum of one far-away factor, and integrate the other.
This contributes \(O_{J,\chi}(H^{-J})\).  The only new jumps are at
the two endpoints of \(I_{\ell,Q}\); both factors there have
normalized distance \(\gg H\), so these point contributions obey the
same bound.  This proves \eqref{eq:scaled-BV-projected}.
\hfill\(\square\)\par
\medskip

We also record the interchange that leads to the Fourier coefficients
below.  The original symmetric principal value agrees with its Abel
regularization.  Indeed, once \(|q|\) is large, the cutoff factor is
one, and pairing \(q\) with \(-q\) gives
\[
 V(Np\alpha)
 \left(\frac1{p\alpha-q}+\frac1{p\alpha+q}\right)
 =O_p(q^{-2}).
\]
Thus the paired tails are absolutely summable, and Abel's theorem
applies to the resulting convergent symmetric series.  To compare the
literal Abel weights with a symmetrically weighted pairing, note that,
for \(q>|p\alpha|\), the two weights are
\(e^{-\lambda_{\rm A}q}e^{\pm\lambda_{\rm A}p\alpha}\).  Replacing both
by \(e^{-\lambda_{\rm A}q}\) changes the paired tail by
\[
 O_p\!\left(
 \lambda_{\rm A}\sum_{q\ge1}\frac{e^{-\lambda_{\rm A}q}}q
 \right)
 =O_p\!\left(\lambda_{\rm A}\log(2/\lambda_{\rm A})\right)=o(1).
\]
The symmetrically weighted paired series is dominated by its
\(O_p(q^{-2})\) tail and therefore converges to the original symmetric
principal value.  We may therefore
first truncate each \(q\)-sum symmetrically and insert the Abel factor
\(\exp(-\lambda_{\rm A}|p\alpha-q|)\), \(\lambda_{\rm A}>0\).  The sum is then finite and,
after the truncation is removed, absolutely integrable.  The Fourier
series \eqref{eq:Vfourier} is absolutely and uniformly convergent, so
it may be inserted term by term before unfolding the residue classes
modulo \(p\).  For the \(\ell\)-th term, the changes of variables
\(x=p\alpha-q\) and then \(v=px\) give the damped transform
\[
 \int_{\R}\frac{\chi(v)}v
 e\left(-\frac{n-\ell Np}{p^2}v\right)
 e^{-\lambda_{\rm A}|v|/p}\dd v.
\]
Letting \(\lambda_{\rm A}\downarrow0\) symmetrically gives
\(R_\chi((n-\ell Np)/p^2)\).  This passage is pointwise in the real
frequency and uniformly bounded; uniform convergence across frequency
zero is neither true nor needed.  More precisely, for the principal
kernel,
\[
 \PV\int_{\R}\frac{e(-sv)e^{-\lambda_{\rm A}|v|/p}}v\dd v
 =-2\ii\arctan\left(\frac{2\pi ps}{\lambda_{\rm A}}\right),
\]
which is uniformly bounded and converges pointwise to
\(-\ii\pi\sgn(s)\).  For the correction, pairing \(v\) and \(-v\) gives
\[
 -2\ii\int_0^C\kappa_\chi(v)
       \frac{\sin(2\pi sv)}v e^{-\lambda_{\rm A}v/p}\dd v,
\]
which is uniformly bounded by Dirichlet's estimate and converges for
each \(s\) as \(\lambda_{\rm A}\downarrow0\).  Since \(p\le N\) is a finite sum and
\(\sum_\ell|v_\ell|<\infty\), dominated convergence now removes the
Abel factor and exchanges the principal value, the Fourier series, and
the cell unfolding.  This is the same symmetric truncation convention
as in Lemma~\ref{lem:transform}.

Using \eqref{eq:Vfourier}, the \(n\)-th Fourier coefficient of the
\(\ell\)-th mark mode of \(Z_N^\chi\), without the factor \(v_\ell\),
is
\begin{equation}\label{eq:B-ell}
 B_\ell(n)=
 \sum_{p\le N}\frac{c_p(-n)}{p^2}
 R_\chi\left(\frac{n-\ell Np}{p^2}\right).
\end{equation}

\smallskip
\noindent\emph{The unshifted mode.}
For \(n\asymp K\), the transition in \eqref{eq:B-ell} with
\(\ell=0\) is \(Q_K=\sqrt K\).  From
\eqref{eq:Rchi-Schwartz}, for a suitable fixed \(C_\chi\ge1\) and
every \(j\),
\[
 |R_\chi(n/p^2)|+
 p\left|\frac{\dd}{\dd p}R_\chi(n/p^2)\right|
 \ll_{j,\chi}
 \begin{cases}
 (p/Q_K)^j,&p<Q_K/C_\chi,\\
 1,&p\ge Q_K/C_\chi.
 \end{cases}
\]
We now spell out the vector-valued Abel summation, including its
boundaries.  Put
\[
 \mathcal H_K=\ell^2\{n\in\Z:K<n\le2K\},\qquad
 u_p=\left(\frac{c_p(-n)}{p^2}\right)_{n\asymp K},\qquad
 w_x=\bigl(R_\chi(n/x^2)\bigr)_{n\asymp K},
\]
and, for a finite upper endpoint \(U\), set
\[
 \mathcal R_{x,U}=\sum_{x<p\le U}u_p.
\]
The exact identity \eqref{eq:near-vector-Abel}, applied with this
multiplier, reads
\begin{equation}\label{eq:Rchi-vector-Abel}
 \sum_{2<p\le U}u_p\odot w_p
 =\mathcal R_{2,U}\odot w_2+
 \int_2^U\mathcal R_{x,U}\odot\frac{\dd w_x}{\dd x}\dd x.
\end{equation}
There is no omitted upper boundary, because
\(\mathcal R_{U,U}=0\).

For all sufficiently large \(K\), put
\(T=K/(\log K)^{B_{\rm tail}}\ge2\) and \(U_0=\min\{N,T\}\); the bounded set
of smaller blocks will be absorbed into the final constant.  The finite-tail
form of \eqref{eq:ram-tail-two-term} gives, uniformly for
\(2\le x\le U_0\),
\begin{equation}\label{eq:Rchi-low-finite-tail}
 \norm{\mathcal R_{x,U_0}}_{\mathcal H_K}
 \ll\frac{\sqrt K}{x}+E_K.
\end{equation}
For a fixed integer \(j>2\), define
\[
 \omega_{K,j}(x)=
 \min\left\{1,\left(\frac{C_\chi x}{Q_K}\right)^j\right\}.
\]
The preceding multiplier bound implies
\begin{equation}\label{eq:Rchi-weight-envelope}
 \norm{w_x}_{\ell^\infty(n\asymp K)}+
 x\norm{\frac{\dd w_x}{\dd x}}_{\ell^\infty(n\asymp K)}
 \ll_{j,\chi}\omega_{K,j}(x).
\end{equation}
The half-line Schwartz bounds also imply the vector-valued total
variation estimate
\begin{equation}\label{eq:Rchi-vector-variation}
 \int_0^\infty\sup_{K\le n\le2K}
 \left|\frac{\dd}{\dd x}R_\chi(n/x^2)\right|\dd x\ll_\chi1.
\end{equation}
To see this, put \(n=rK\), \(1\le r\le2\), and \(t=K/x^2\); the
left side is bounded by a constant times
\(\int_0^\infty\sup_{1\le r\le2}r|R_\chi'(rt)|\dd t\).
This is finite by \eqref{eq:Rchi-Schwartz}.  Applying
\eqref{eq:Rchi-vector-Abel}, \eqref{eq:Rchi-low-finite-tail}, and
\eqref{eq:Rchi-weight-envelope}, and using
\[
 \frac{\omega_{K,j}(2)}2+
 \int_2^\infty\frac{\omega_{K,j}(x)}{x^2}\dd x
 \ll_{j,\chi}\frac1{Q_K},
\]
we obtain
\begin{equation}\label{eq:Rchi-low-Abel-bound}
 \left\|\sum_{2<p\le U_0}u_p\odot w_p\right\|_{\mathcal H_K}
 \ll_{j,\chi}\frac{\sqrt K}{Q_K}+E_K
 =O_\chi(1)+o(1).
\end{equation}
The \(E_K\)-part includes the boundary at \(2\) and is bounded using
\eqref{eq:Rchi-vector-variation}.

If \(T<N\), apply the same exact identity to \((T,N]\).  The finite
form of \eqref{eq:ram-tail-terminal} gives
\[
 \norm{\mathcal R_{x,N}}_{\mathcal H_K}
 \ll\frac{\sqrt K(\log K)^{3/2}}x,
 \qquad T\le x\le N,
\]
where now \(\mathcal R_{x,N}=\sum_{x<p\le N}u_p\).  Moreover,
\[
 \left|\frac{\dd}{\dd x}R_\chi(n/x^2)\right|
 \ll_\chi\min(x^{-1},Kx^{-3}).
\]
The lower boundary at \(T\) is
\(O(\sqrt K(\log K)^{3/2}/T)=O(E_K)\), since \(R_\chi\) is bounded;
the upper boundary again vanishes.  The remaining Abel integral is
\[
 \ll \sqrt K(\log K)^{3/2}K\int_T^\infty x^{-4}\dd x
 =o(E_K).
\]
The finitely many terms \(p\le2\) satisfy an \(O_\chi(1)\) bound by
the first multiplier estimate.  We have therefore proved
\begin{equation}\label{eq:B0-block}
 \norm{B_0}_{\ell^2(n\asymp K)}\ll_\chi1.
\end{equation}
We also record the high-frequency tail rather than subsume it in
\eqref{eq:B0-block}.  Suppose that \(Q_K=\sqrt K>N\), equivalently
\(K>N^2\).  For all sufficiently large \(N\), this implies \(T>N\):
the function \(x/(\log x)^{B_{\rm tail}}\) is increasing in this range, and at
\(x=N^2\) it is larger than \(N\).
Furthermore, once \(N\) is sufficiently large, the function
\(x\mapsto(\log x)^{B_{\rm tail}+3/2}x^{-1/2}\) is decreasing on
\([N^2,\infty)\).  Hence
\[
 \sup_{K\ge N^2}E_K
 \le \frac{(\log N^2)^{B_{\rm tail}+3/2}}N=o(1),
\]
so every use of \(E_K=o(1)\) below is uniform over the high-frequency
blocks.  Apply
\eqref{eq:Rchi-vector-Abel} with \(U=N\).  Every \(x\le N\) lies below
\(Q_K\), so \eqref{eq:Rchi-weight-envelope} and
\eqref{eq:Rchi-low-finite-tail} give, for each fixed \(j>2\),
\begin{align}
 \left\|\sum_{2<p\le N}u_p\odot w_p\right\|_{\mathcal H_K}
 &\ll_{j,\chi}
 \sqrt K Q_K^{-j}
 \left(1+\int_2^Nx^{j-2}\dd x\right)
 +E_K\left(\frac N{Q_K}\right)^j\notag\\
 &\ll_{j,\chi}
 \left(\frac N{Q_K}\right)^{j-1}.
 \label{eq:Rchi-high-frequency-block}
\end{align}
In the last step we used \(\sqrt K/Q_K=1\),
\(E_K=o(1)\), and \(N/Q_K<1\).  The terms \(p\le2\) obey the same
bound directly from \eqref{eq:Rchi-Schwartz}, after increasing the
fixed decay exponent if necessary.  On the dyadic blocks
\(K_r=2^rN^2\), the square of
\eqref{eq:Rchi-high-frequency-block} is
\(O_{j,\chi}(2^{-r(j-1)})\), and is therefore summable.  The range
\(K\le N^2\) contains \(O(\log N)=O(L)\) dyadic blocks, each bounded
by \eqref{eq:B0-block}.

Finally, the reality of \(\chi(v)/v\) gives
\(R_\chi(-t)=\overline{R_\chi(t)}\), and Ramanujan sums are real and
even.  Hence \(B_0(-n)=\overline{B_0(n)}\), so the negative-frequency
blocks have the same bounds.  At frequency zero,
\(B_0(0)=0\) by \eqref{eq:Rchi-jump}.  Combining the positive,
negative, and zero frequencies proves
\begin{equation}\label{eq:B0-total}
 \sum_n|B_0(n)|^2\ll_\chi L.
\end{equation}

\smallskip
\noindent\emph{Shifted modes.}
Fix \(\ell\ne0\).  For \(p,p'\asymp Q\), apply the auxiliary
scaled-multiplier estimate above with \(a=\ell Np\) and
\(a'=\ell Np'\).  It gives, for
\[
 W_{p,p'}(t)=
 R_\chi\left(\frac{t-\ell Np}{p^2}\right)
 \overline{R_\chi\left(\frac{t-\ell Np'}{p'^2}\right)},
\]
\begin{equation}\label{eq:BV-product}
 \norm{W_{p,p'}}_\infty+\TV_{\R}(W_{p,p'})
 \ll_{J,\chi}
 \left(1+\frac{|\ell|N|p-p'|}{Q^2}\right)^{-J}.
\end{equation}
Before making the dyadic decomposition, remove the singleton \(p=1\).
Since \(c_1(n)=1\), its contribution is
\[
 B_{\ell}^{[1]}(n)=R_\chi(n-\ell N).
\]
Translation by the integer \(\ell N\), together with
\eqref{eq:Rchi-Schwartz} and \(R_\chi(0)=0\), gives the uniform bound
\begin{equation}\label{eq:shifted-p-one}
 \sum_{n\in\Z}|B_{\ell}^{[1]}(n)|^2
 =\sum_{m\in\Z}|R_\chi(m)|^2\ll_\chi1.
\end{equation}
Thus it remains to decompose the terms \(2\le p\le N\) into the
blocks \(Q<p\le2Q\) below; the singleton can be restored at the end
by the triangle inequality.
The function \(W_{p,p'}\) tends to zero at both infinities.  Therefore
Lemmas~\ref{lem:discrete-BV} and~\ref{lem:incomplete}, together with
\eqref{eq:BV-product}, imply, when \(p\ne p'\),
\begin{equation}\label{eq:shifted-cross}
 \left|\sum_{n\in\Z}c_p(n)c_{p'}(n)W_{p,p'}(n)\right|
 \ll_{J,\chi}\sigma_1(p)\sigma_1(p')
 \left(1+\frac{|\ell|N|p-p'|}{Q^2}\right)^{-J}.
\end{equation}
Let \(B_{\ell,Q}\) denote the dyadic \(Q<p\le2Q\) part of
\eqref{eq:B-ell}.  We first write out its diagonal contribution.
Put \(a=\ell Np\).  Since \(a/p=\ell N\in\Z\), every interval
\((a+kp,a+(k+1)p]\cap\Z\) is a complete residue system modulo \(p\),
and hence
\[
 \sum_{a+kp<n\le a+(k+1)p}c_p(n)^2=p\varphi(p).
\]
The half-line bounds \eqref{eq:Rchi-Schwartz} now give
\begin{align}
 \sum_{n\in\Z}c_p(n)^2
 \left|R_\chi\left(\frac{n-\ell Np}{p^2}\right)\right|^2
 &\le p\varphi(p)\sum_{k\in\Z}
 \sup_{k/p<x\le(k+1)/p}|R_\chi(x)|^2\notag\\
 &\ll_\chi p^2\varphi(p).
 \label{eq:shifted-diagonal-periods}
\end{align}
Consequently the diagonal part of \(\norm{B_{\ell,Q}}_2^2\) is
\[
 \ll_\chi\sum_{Q<p\le2Q}\frac{\varphi(p)}{p^2}\ll_\chi1.
\]

For the off-diagonal, put \(\lambda=|\ell|N/Q^2\), use
\(\sigma_1(p)/p\ll\log\log(3Q)\), and choose \(J>2\).  Equation
\eqref{eq:shifted-cross} gives
\begin{align}
 &\sum_{\substack{p,p'\asymp Q\\p\ne p'}}
 \frac{\sigma_1(p)\sigma_1(p')}{p^2p'^2}
 (1+\lambda|p-p'|)^{-J}\notag\\
 &\qquad\ll
 \frac{(\log\log(3Q))^2}{Q^2}
 Q\sum_{1\le h\le Q}(1+\lambda h)^{-J}\notag\\
 &\qquad\ll
 (\log\log(3Q))^2
 \min\left(1,\frac{Q}{|\ell|N}\right).
 \label{eq:shifted-offdiagonal-sum}
\end{align}
Combining this with \eqref{eq:shifted-diagonal-periods} proves, uniformly
for \(\ell\ne0\), the slightly stronger form
\begin{equation}\label{eq:shifted-Q}
 \norm{B_{\ell,Q}}_2^2
 \ll_\chi
 1+(\log\log(3Q))^2
 \min\left(1,\frac{Q}{|\ell|N}\right).
\end{equation}

It remains to square-sum in \(Q\).  Let
\[
 Q_0=\frac{N}{L^{C_{\rm far}}},\qquad C_{\rm far}>10,
\]
and let \(\Pi_{\ell,Q}\) project onto
\[
 \sgn(n)=\sgn(\ell),\qquad
 \frac{|\ell|NQ}{2}\le|n|\le3|\ell|NQ.
\]
Put \(H=|\ell|N/Q\) and let \(E_{\ell,Q}\) be the complement in
\(\R\) of the displayed carrier annulus.  Expanding by Plancherel gives
the exact identity
\begin{equation}\label{eq:carrier-expansion}
 \begin{split}
 \norm{(1-\Pi_{\ell,Q})B_{\ell,Q}}_2^2
 ={}&\sum_{p,p'\asymp Q}\frac1{p^2p'^2}
 \sum_{n\in\Z}c_p(n)c_{p'}(n)\\
 &\hspace{18mm}\cdot
 \1_{E_{\ell,Q}}(n)W_{p,p'}(n).
 \end{split}
\end{equation}
For \(p\ne p'\), the projected estimate
\eqref{eq:scaled-BV-projected} and Lemmas~\ref{lem:discrete-BV}
and~\ref{lem:incomplete} bound the inner sum by
\begin{equation}\label{eq:carrier-offdiagonal-inner}
 O_{J,\chi}\bigl(\sigma_1(p)\sigma_1(p')H^{-J}\bigr).
\end{equation}

For \(p=p'\), every retained \(n\) satisfies
\(|n-\ell Np|/p^2\gg H\).  Subdivide into the periods used in
\eqref{eq:shifted-diagonal-periods} and take a Schwartz exponent larger
than \(J+1\).  The remaining supremum sum is bounded by
\[
 p\int_{|x|\gg H}(1+|x|)^{-J-2}\dd x+O(H^{-J-2})
 \ll_JpH^{-J}.
\]
It follows that
\begin{equation}\label{eq:carrier-diagonal-inner}
 \sum_{n\in E_{\ell,Q}\cap\Z}c_p(n)^2
 \left|R_\chi\left(\frac{n-\ell Np}{p^2}\right)\right|^2
 \ll_{J,\chi}p^2\varphi(p)H^{-J}.
\end{equation}
After applying the outside weights in
\eqref{eq:carrier-expansion}, the diagonal contribution is
\(O_{J,\chi}(H^{-J})\), because
\(\sum_{p\asymp Q}\varphi(p)/p^2\ll1\).  The off-diagonal contribution
from \eqref{eq:carrier-offdiagonal-inner} is
\(O_{J,\chi}((\log\log(3Q))^2H^{-J})\).  We have therefore proved,
including both projection-boundary jumps, that for every \(J\),
\begin{equation}\label{eq:carrier-leak}
 \norm{(1-\Pi_{\ell,Q})B_{\ell,Q}}_2^2
 \ll_{J,\chi}
 (1+\log\log(3Q))^2
 \left(\frac{Q}{|\ell|N}\right)^J.
\end{equation}

The carrier annuli have bounded overlap as \(Q\) ranges dyadically:
for each integer frequency \(n\), at most a fixed number of dyadic
values of \(Q\) can satisfy \(n\in I_{\ell,Q}\), since those values
have ratio at most \(6\).  Therefore
\begin{equation}\label{eq:carrier-bounded-overlap}
 \begin{split}
 \left\|\sum_Q\Pi_{\ell,Q}B_{\ell,Q}\right\|_2^2
 &=\sum_n\left|\sum_Q
 \1_{I_{\ell,Q}}(n)B_{\ell,Q}(n)\right|^2\\
 &\ll\sum_Q\norm{B_{\ell,Q}}_2^2.
 \end{split}
\end{equation}
For \(Q\le Q_0\), the sum of \eqref{eq:shifted-Q} over dyadic \(Q\) is
\(O_\chi(L)\): the constant terms contribute \(O(L)\), while
\[
 \sum_{Q\le Q_0}(\log\log(3Q))^2\frac{Q}{|\ell|N}=o(1).
\]
Moreover, after taking square roots in \eqref{eq:carrier-leak},
\[
 \sum_{Q\le Q_0}
 \norm{(1-\Pi_{\ell,Q})B_{\ell,Q}}_2=o(1)
\]
on choosing \(J\) large.  Bounded overlap of the projected annuli
in \eqref{eq:carrier-bounded-overlap} therefore gives
\begin{equation}\label{eq:low-Q}
 \left\|\sum_{Q\le Q_0}B_{\ell,Q}\right\|_2^2\ll_\chi L.
\end{equation}
There are \(O(\log L)\) blocks above \(Q_0\).  Each has squared norm
\(O_\chi((\log L)^2)\), so the triangle inequality gives
\[
 \sup_{\ell\ne0}\sum_n|B_\ell(n)|^2
 \ll_\chi L+(\log L)^4=o(L^2).
\]
Minkowski, \(\sum_{\ell\ne0}|v_\ell|<\infty\), and
\eqref{eq:B0-total} prove
\[
 \norm{Z_N^\chi}_2^2\ll_\chi L+(\log L)^4=o(L^2).
\]
Together with \eqref{eq:far-window-L2}, the same holds for
\(Y_N^\chi\).

Finally remove the smoothing.  The difference between \(Y_N^\chi\)
and the sharp sum in \eqref{eq:far} is supported on
\[
 \varepsilon-\delta<|p\delta_p|<\varepsilon.
\]
Every active \(p\) is a convergent denominator, so at most \(O(L)\)
are active at one \(\alpha\), since \(Q_n\ge F_{n+1}\) for the Fibonacci
numbers \(F_n\).  Each summand is \(O_\varepsilon(1)\), and
\[
 \int_0^1\sum_{p\le N}
 \1_{\{\varepsilon-\delta<|p\delta_p|<\varepsilon,\,
       (p,q_p)=1\}}\dd\alpha
 \le2\delta\sum_{p\le N}\frac{\varphi(p)}{p^2}
 \ll\delta L.
\]
Pointwise Cauchy--Schwarz thus bounds the squared \(L^2\)-norm of the
difference by \(O_\varepsilon(\delta L^2)\).  First let \(N\to\infty\)
and then \(\delta\downarrow0\).  This proves \eqref{eq:far}.

It remains to prove the quantified uniformity
\eqref{eq:far-uniform}.  Fix
\(0<\delta<\varepsilon_0/2\).  For every
\(t\in[\varepsilon_0,\varepsilon_1]\), construct
\(\chi_{t,\delta}\) from one fixed smooth transition profile so that it
vanishes on \(|v|\le t-\delta\) and equals one on \(|v|\ge t\).
All these profiles vanish on a fixed neighborhood of zero, and their
transition points remain in a fixed compact set.  Hence, for each fixed
\(r,J\),
\begin{equation}\label{eq:uniform-chi-seminorms}
 \sup_{t\in[\varepsilon_0,\varepsilon_1]}
 C_{r,J,\chi_{t,\delta}}<\infty,
\end{equation}
where the bound may depend on the presently fixed \(\delta\).  Every
periodization and multiplier estimate above is consequently uniform in
\(t\).  In particular, \eqref{eq:far-window-L2} and the bound for
\(Z_N^\chi\) give the explicit common estimate
\[
 \sup_{t\in[\varepsilon_0,\varepsilon_1]}
 \norm{Y_N^{\chi_{t,\delta}}}_2^2
 \ll L\log L+C_\delta\bigl(L+(\log L)^4\bigr)=o_\delta(L^2).
\]
Equivalently,
\begin{equation}\label{eq:uniform-smoothed-far}
 \sup_{t\in[\varepsilon_0,\varepsilon_1]}
 \norm{Y_N^{\chi_{t,\delta}}}_2=o_\delta(L).
\end{equation}

Let \(F_N(t)\) denote the sharp sum in \eqref{eq:far-uniform}, and put
\(D_{N,t}=F_N(t)-Y_N^{\chi_{t,\delta}}\).  This difference is supported
on
\[
 t-\delta<|p\delta_p|<t.
\]
Each active summand is \(O_{\varepsilon_0}(1)\).  Moreover, its reduced
fraction satisfies
\[
 \left|\alpha-\frac{q_p}{p}\right|
 <\frac{\varepsilon_1}{p^2}<\frac1{2p^2}.
\]
Legendre's criterion therefore makes \(q_p/p\) a continued-fraction
convergent.  Since consecutive convergent denominators satisfy
\(Q_{n+2}\ge2Q_n\), the number of active \(p\le N\) is \(O(L)\),
uniformly in \(\alpha\) and \(t\).  On the other hand, direct cell
integration gives, uniformly in \(t\),
\begin{equation}\label{eq:uniform-transition-first-moment}
 \int_0^1\sum_{p\le N}
 \1_{\{t-\delta<|p\delta_p|<t,\ (p,q_p)=1\}}\dd\alpha
 \le2\delta\sum_{p\le N}\frac{\varphi(p)}{p^2}+O(\delta)
 \ll\delta L.
\end{equation}
The \(O(\delta)\) term harmlessly covers the endpoint cells for
\(p=1\).  Pointwise Cauchy--Schwarz, followed by
\eqref{eq:uniform-transition-first-moment}, yields
\begin{equation}\label{eq:uniform-transition-L2}
 \sup_{t\in[\varepsilon_0,\varepsilon_1]}
 \norm{D_{N,t}}_2^2\ll_{\varepsilon_0}\delta L^2.
\end{equation}
Combining \eqref{eq:uniform-smoothed-far} and
\eqref{eq:uniform-transition-L2},
\[
 \limsup_{N\to\infty}
 \sup_{t\in[\varepsilon_0,\varepsilon_1]}
 L^{-1}\norm{F_N(t)}_2
 \ll_{\varepsilon_0}\sqrt\delta.
\]
Letting \(\delta\downarrow0\) proves \eqref{eq:far-uniform}.  Thus the
order of limits is: first \(N\to\infty\) with \(\delta\) fixed, and
then \(\delta\downarrow0\).
\end{proof}

\subsection{The complete small-jump estimate}

\begin{proposition}[Minor-arc truncation]\label{prop:minor}
For every \(\gamma>0\),
\begin{equation}\label{eq:minor-final}
 \lim_{A\to\infty}\limsup_{N\to\infty}
 \Leb\left\{\alpha:
 \frac1L\left|
 \sum_{\substack{p\le N,\ (p,q_p)=1\\L|p\delta_p|>A}}
 \frac{V(N\delta_p)}{p\delta_p}\right|>\gamma\right\}=0.
\end{equation}
\end{proposition}

\begin{proof}
Fix \(0<\varepsilon<1/4\), and abbreviate
\[
 v_p=p\delta_p,\qquad \xi_p=Lv_p,\qquad
 u_p=N\delta_p\bmod1.
\]
Let
\[
 F_N(t)=
 \sum_{\substack{p\le N,\ (p,q_p)=1\\|v_p|\ge t}}
 \frac{V(u_p)}{v_p}.
\]
Proposition~\ref{prop:far} gives
\(\|F_N(\varepsilon/2)\|_2=o(L)\).

For all sufficiently large \(N\), the fixed-profile cutoff used in
Proposition~\ref{prop:near} yields the exact decomposition, outside null
boundary sets,
\begin{equation}\label{eq:minor-decomposition}
 \sum_{\substack{p\le N,\ (p,q_p)=1\\|\xi_p|>A}}
 \frac{V(u_p)}{v_p}
 =
 F_N(\varepsilon/2)
 +Y^{\mathrm{near}}_{N,A,\varepsilon}
 +E^{\mathrm{lo}}_{N,A}-E^{\mathrm{up}}_{N,\varepsilon},
\end{equation}
where
\begin{align*}
 E^{\mathrm{lo}}_{N,A}
 &=
 \sum_{\substack{p\le N,\ (p,q_p)=1\\A<|\xi_p|<2A}}
 \left(1-\eta(|\xi_p|/A)\right)\frac{V(u_p)}{v_p},\\
 E^{\mathrm{up}}_{N,\varepsilon}
 &=
 \sum_{\substack{p\le N,\ (p,q_p)=1\\
                  \varepsilon/2<|v_p|<\varepsilon}}
 \beta(|v_p|/\varepsilon)\frac{V(u_p)}{v_p}.
\end{align*}
This identity, rather than a containment argument, is needed because
cancellation does not compare a weighted transition sum with its
containing sharp band.

The near term satisfies
\[
 \limsup_{N\to\infty}
 L^{-2}\|Y^{\mathrm{near}}_{N,A,\varepsilon}\|_2^2
 \ll_\varepsilon A^{-1}
\]
by Proposition~\ref{prop:near}.  For the lower transition, define
\(\mathcal P\) to be the limiting Poisson process in
Lemma~\ref{lem:marked-poisson}, and put
\[
 h_A(\xi,u)=
 \left(1-\eta(|\xi|/A)\right)
 \1_{\{A<|\xi|<2A\}}\frac{V(u)}{\xi}.
\]
Lemma~\ref{lem:marked-poisson} and the continuous mapping theorem
(the boundary has zero limiting intensity) give
\[
 \frac{E^{\mathrm{lo}}_{N,A}}L
 \Longrightarrow \int h_A\,\dd\mathcal P.
\]
The process in Lemma~\ref{lem:marked-poisson} uses \(p<N\); the omitted
\(p=N\) term is nonzero on a set of measure \(O_A((LN)^{-1})\), so it
does not affect this convergence.
The limiting shot has mean zero by symmetry in \(\xi\), and
\[
 \Var\!\left(\int h_A\,\dd\mathcal P\right)
 =\frac1{\zeta(2)}\int h_A(\xi,u)^2\dd\xi\dd u
 \ll\frac1A.
\]
Thus Portmanteau and Chebyshev bound its contribution to
\eqref{eq:minor-final} by \(O(A^{-1}\gamma^{-2})\).

It remains to justify the upper transition.  Put
\(r_0=\varepsilon/2\), \(r_1=\varepsilon\), and
\(h(t)=\beta(t/\varepsilon)\).  Stieltjes layer cake gives exactly
\[
 E^{\mathrm{up}}_{N,\varepsilon}
 =h(r_0)F_N(r_0)+\int_{r_0}^{r_1}h'(t)F_N(t)\dd t.
\]
This identity follows first pointwise in \(\alpha\), by applying the
one-variable Stieltjes identity to each of the finitely many shots and
then summing.  As a map from \([r_0,r_1]\) to the separable space
\(L^2(\T)\), \(t\mapsto F_N(t)\) is strongly measurable: its joint
representative is measurable in \((t,\alpha)\).  Moreover, for fixed
\(N\), throughout this interval
\[
 |F_N(t,\alpha)|\le \frac{N\|V\|_\infty}{r_0},
\]
so it is Bochner integrable.  The pointwise identity is therefore also
an identity in \(L^2(\T)\).
Minkowski's inequality and \eqref{eq:far-uniform} therefore give
\[
 \frac{\norm{E^{\mathrm{up}}_{N,\varepsilon}}_2}{L}
 \le
 \left(|h(r_0)|+\int_{r_0}^{r_1}|h'(t)|\dd t\right)
 \sup_{r_0\le t\le r_1}\frac{\norm{F_N(t)}_2}{L}
 \longrightarrow0.
\]
Apply Chebyshev to the near term, use the preceding bounds in
\eqref{eq:minor-decomposition}, and finally let \(A\to\infty\).
This proves \eqref{eq:minor-final}.
\end{proof}

\section{The marked Poisson process and the limit law}\label{sec:poisson}

\subsection{A continued-fraction marked Poisson theorem}

Let \(m\) be normalized Haar measure on \(\T\).  We prove the
one-dimensional marked theorem directly; no rank-two specialization of a
Rogers formula is used.

\begin{lemma}[Signed marked resonances]\label{lem:marked-poisson}
For \(\alpha\) uniform on \(\T\), the point process
\begin{equation}\label{eq:point-process-time}
 \mathcal Q_N=
 \sum_{\substack{1\le p<N\\(p,q_p)=1}}
 \delta_{\left(\log p/L,\;Lp\delta_p,\;N\delta_p\bmod1\right)}
\end{equation}
converges vaguely on
\([0,1]\times(\R\setminus\{0\})\times\T\) to a Poisson process with
intensity
\begin{equation}\label{eq:PPP-intensity-time}
 \lambda\,dt\,d\xi\,dm(u),\qquad
 \lambda=\frac6{\pi^2}=\frac1{\zeta(2)}.
\end{equation}
Consequently,
\begin{equation}\label{eq:point-process}
 \mathcal P_N=
 \sum_{\substack{1\le p<N\\(p,q_p)=1}}
 \delta_{\left(Lp\delta_p,\;N\delta_p\bmod1\right)}
\end{equation}
converges vaguely on
\((\R\setminus\{0\})\times\T\) to the Poisson process with intensity
\begin{equation}\label{eq:PPP-intensity}
 \lambda\,d\xi\,dm(u).
\end{equation}
\end{lemma}

\begin{proof}
\emph{Roadmap.}
We first establish the unmarked factorial measures in deterministic
continued-fraction index time and transfer them from Gauss to Lebesgue
measure.  We then prove asymptotic Haar marking by an oscillatory
cylinder estimate, treating early and late configurations separately.
Finally, factorial measures are upgraded to point-process convergence,
index time is replaced by denominator time, and the endpoint time layers
are restored by a direct rational-cell estimate.

\medskip
\emph{Convergents and rare digits.}
Write
\[
 \alpha=[0;a_1,a_2,\ldots]
\]
and denote its regular convergents by \(P_n/Q_n\), with
\[
 P_{-1}=1,\quad P_0=0,\qquad Q_{-1}=0,\quad Q_0=1.
\]
Let \(T(x)=\{1/x\}\), \(x_n=T^n\alpha\), and
\(y_n=Q_{n-1}/Q_n\).  The determinant identities for convergents give
\begin{equation}\label{eq:cf-error}
 Q_n\alpha-P_n=(-1)^n\frac{\theta_n}{Q_n},\qquad
 \theta_n=\frac{x_n}{1+x_ny_n}
 =\frac1{a_{n+1}+x_{n+1}+y_n}.
\end{equation}
These identities include \(n=0\), with \(y_0=0\).

Fix \(B<\infty\).  If a point in \eqref{eq:point-process} satisfies
\(|Lp\delta_p|\le B\), then, for all sufficiently large \(N\),
\[
 \left|\alpha-\frac{q_p}{p}\right|
 <\frac1{2p^2}.
\]
Since \(q_p/p\) is reduced, Legendre's criterion
\cite[Chapter~II, \S7, Theorem~19, p.~30]{Khinchin} says that it is a
convergent.  Thus, on
every fixed \(\xi\)-window, \eqref{eq:point-process-time} is exactly
\begin{equation}\label{eq:cf-process}
 \sum_{Q_n<N}
 \delta_{\left(\log Q_n/L,\;(-1)^nL\theta_n,\;
              (-1)^nN\theta_n/Q_n\bmod1\right)}.
\end{equation}

Let \(\nu_G\) be Gauss measure,
\[
 d\nu_G(x)=\frac{dx}{\log2\,(1+x)}.
\]
Set
\[
 \kappa=\frac{12\log2}{\pi^2}.
\]
The continued-fraction digits are exponentially \(\psi\)-mixing: there are
\(C>0\) and \(\vartheta>1\) such that, if
\[
 F\in\sigma(a_1,\ldots,a_j),\qquad
 H\in\sigma(a_{j+d},a_{j+d+1},\ldots),
\]
then
\begin{equation}\label{eq:gauss-mixing}
 |\nu_G(F\cap H)-\nu_G(F)\nu_G(H)|
 \le C\vartheta^{-d}\nu_G(F)\nu_G(H).
\end{equation}
We use the form and notation recorded in
\cite[(1.7)]{GhoshKirsebomRoy}.  We record two consequences of
\eqref{eq:gauss-mixing} that will be used repeatedly.  Suppose that
\(\mathscr A\) and \(\mathscr H\) are \(\sigma\)-fields for which
\[
 \left|\nu_G(A\cap B)-\nu_G(A)\nu_G(B)\right|
 \le \eta\,\nu_G(A)\nu_G(B)
 \qquad(A\in\mathscr A,\ B\in\mathscr H).
\]
Then, for arbitrary complex-valued
\(U_1\in L^1(\mathscr A)\) and \(U_2\in L^1(\mathscr H)\),
\begin{equation}\label{eq:functional-psi-mixing}
 \left|\E_{\nu_G}(U_1U_2)
       -(\E_{\nu_G}U_1)(\E_{\nu_G}U_2)\right|
 \le 2\eta\,\E_{\nu_G}|U_1|\,\E_{\nu_G}|U_2|.
\end{equation}
Indeed, for nonnegative simple functions this follows by expanding both
functions into indicator functions and applying the event inequality
term by term.  Monotone convergence gives the result for arbitrary
nonnegative integrable functions.  For a complex function \(U_1\), write
\[
 U_1=(\Re U_1)^+-(\Re U_1)^-
   +\ii\bigl((\Im U_1)^+-(\Im U_1)^-\bigr),
\]
and similarly for \(U_2\).  The sums of the \(L^1\)-norms of the
four nonnegative parts are at most \(\sqrt2\,\E|U_1|\) and
\(\sqrt2\,\E|U_2|\), respectively, which proves
\eqref{eq:functional-psi-mixing}.

A second consequence is the following quasi-Bernoulli bound.  If
\(m_1<\cdots<m_s\) and \(E_j\in\sigma(a_{m_j})\), then
\begin{equation}\label{eq:digit-quasi-bernoulli}
 \nu_G\left(\bigcap_{j=1}^sE_j\right)
 \le C_s\prod_{j=1}^s\nu_G(E_j).
\end{equation}
To see this, expose the events from left to right and apply
\eqref{eq:gauss-mixing} at each step.  Every gap is at least one, so
each application costs at most \(1+C\vartheta^{-1}\).  This explains
precisely what is meant below by the quasi-Bernoulli bound.

For a compact interval \(J\Subset(0,\infty)\), the last formula in
\eqref{eq:cf-error} gives, on the event \(L\theta_n\in J\),
\begin{equation}\label{eq:digit-replacement}
 \left|L\theta_n-\frac{L}{a_{n+1}}\right|
 \ll_J\frac1L.
\end{equation}
The same conclusion, with a constant depending on a fixed enlargement
of \(J\), holds when \(L/a_{n+1}\in J\).  Indeed, either event forces
\(a_{n+1}\asymp_J L\), and
\[
 \left|\frac{L}{a_{n+1}}
       -\frac{L}{a_{n+1}+x_{n+1}+y_n}\right|
 \le\frac{2L}{a_{n+1}^2}.
\]
Thus, whenever either of the two indicators differs from the other,
both value coordinates lie in one fixed compact subset of
\((0,\infty)\), and one of them lies within \(O_J(L^{-1})\) of an
endpoint of \(J\).
Moreover, direct summation of
\[
 \nu_G(a_1=q)
 =\frac1{\log2}
 \log\frac{(q+1)^2}{q(q+2)}
 =\frac1{\log2}\left(\frac1{q^2}+O(q^{-3})\right)
\]
shows that
\begin{equation}\label{eq:rare-onepoint}
 \nu_G\left\{\frac{L}{a_1}\in J\right\}
 =\frac{|J|}{L\log2}+O_J(L^{-2}).
\end{equation}

\medskip
\noindent\emph{Unmarked factorial measures.}
We give the factorial-moment argument explicitly.  Fix \(r\), compact
index-time intervals \(I_j\), compact value intervals
\(J_j\Subset(0,\infty)\), and prescribed parities.  Put
\(b_L=\lceil C_{\rm mix}\log L\rceil\), with \(C_{\rm mix}\) large, and sum over
distinct positions \(n_j\in\kappa LI_j\) having the prescribed parities.
For the digit events
\[
 A_{n,j}=\{L/a_{n+1}\in J_j\},
\]
\eqref{eq:rare-onepoint} gives
\[
 \nu_G(A_{n,j})=\frac{|J_j|}{L\log2}+O_{J_j}(L^{-2}).
\]
If the ordered positions have successive gaps at least \(b_L\), repeated
use of \eqref{eq:gauss-mixing} gives
\begin{equation}\label{eq:long-tuple-product}
 \nu_G\!\left(\bigcap_{j=1}^rA_{n_j,j}\right)
 =\prod_{j=1}^r\nu_G(A_{n_j,j})
  \left(1+O_r(\vartheta^{-b_L})\right).
\end{equation}
If some gap is smaller than \(b_L\), there are only
\(O_r(b_LL^{r-1})\) possible tuples.  Applying
\eqref{eq:gauss-mixing} successively with gap one bounds every such
intersection by \(O_r(L^{-r})\).  Hence all short-gap tuples contribute
\begin{equation}\label{eq:short-tuples}
 O_r(b_L/L)=o(1).
\end{equation}
For a value-window endpoint \(z>0\), the replacement error in
\eqref{eq:digit-replacement} is contained in
\[
 \{\operatorname{dist}(L/a_{n+1},z)\le C/L\},
\]
whose probability is \(O_z(L^{-2})\), by the same direct digit
summation as in \eqref{eq:rare-onepoint}.  Combining this boundary event
with \(r-1\) ordinary rare events and using the gap-one bound shows that
the sum of all replacement errors is
\(O_r(L^rL^{-(r+1)})=O_r(L^{-1})\).

Thus \eqref{eq:long-tuple-product}, \eqref{eq:short-tuples}, and the
Riemann sums over the index windows prove convergence of every unmarked
factorial measure in deterministic index time.  The number of integers
of either parity in \(\kappa L I_j\) is
\(\frac12\kappa L|I_j|+O(1)\), and
\(\kappa/(2\log2)=6/\pi^2=\lambda\).  Explicitly, for prescribed
parities and labeled intervals as above,
\begin{equation}\label{eq:unmarked-factorial-limit}
 \sum_{\substack{n_j\in\kappa L I_j\\
                  n_1,\ldots,n_r\ {\rm distinct}\\
                  n_j\ {\rm of\ the\ prescribed\ parity}}}
 \nu_G\left(\bigcap_{j=1}^r
       \{L\theta_{n_j}\in J_j\}\right)
 \longrightarrow
 \lambda^r\prod_{j=1}^r|I_j|\,|J_j|,
 \qquad \lambda=\frac6{\pi^2}.
\end{equation}
We now prove explicitly
the uniform moment bounds that will be used below.  Let \(\mathscr I_L\) be
any union of a fixed number of deterministic index-time windows, so that
\(\#\mathscr I_L=O(L)\), let \(K\Subset(0,\infty)\), and put
\[
 R_L=\sum_{n\in\mathscr I_L}\1_{\{L\theta_n\in K\}}.
\]
Since
\[
 \theta_n=\frac1{a_{n+1}+x_{n+1}+y_n},\qquad
 0<x_{n+1}+y_n<2,
\]
the event \(L\theta_n\in K\) implies
\(c_KL\le a_{n+1}\le C_KL\) for suitable positive constants
\(c_K,C_K\).  The latter one-digit event has Gauss probability
\(O_K(L^{-1})\).  Therefore, for every fixed \(s\ge1\),
\eqref{eq:digit-quasi-bernoulli} gives
\begin{align}
 \E_{\nu_G}(R_L)_s
 &=\sum_{\substack{n_1,\ldots,n_s\in\mathscr I_L\\
                  \text{all distinct}}}
   \nu_G\left(\bigcap_{j=1}^s
        \{L\theta_{n_j}\in K\}\right) \notag\\
 &\le
 \sum_{\substack{n_1,\ldots,n_s\in\mathscr I_L\\
                  \text{all distinct}}}
   \nu_G\left(\bigcap_{j=1}^s
        \{c_KL\le a_{n_j+1}\le C_KL\}\right)
 \ll_{K,s}L^sL^{-s}\ll_{K,s}1.
 \label{eq:rare-factorial-bound}
\end{align}
Here \((x)_s=x(x-1)\cdots(x-s+1)\).  Since
\(d\Leb/d\nu_G=\log2\,(1+x)\) is bounded, the same estimate
holds under Lebesgue measure.  Finally,
\begin{equation}\label{eq:ordinary-factorial-identity}
 R_L^s=\sum_{j=0}^s
 \left\{\begin{matrix}s\\ j\end{matrix}\right\}(R_L)_j,
\end{equation}
where the braces are Stirling numbers of the second kind.  We have
therefore proved
\begin{equation}\label{eq:rare-count-moments}
 \sup_L\E_{\nu_G}R_L^s+\sup_L\E_{\rm Leb}R_L^s<\infty
 \qquad(s\ge1).
\end{equation}

\emph{Denominator time and the unmarked intensity.}
The exact identity
\begin{equation}\label{eq:cf-roof}
 \log(Q_n+Q_{n-1}x_n)
 =\sum_{j=0}^{n-1}-\log x_j
\end{equation}
holds.  The Gauss map preserves \(\nu_G\)
\cite[Chapter~III]{Khinchin}, and its ergodicity also follows directly
from \eqref{eq:gauss-mixing}.  Since \(-\log x\in L^1(\nu_G)\),
Birkhoff's ergodic theorem \cite{Birkhoff} implies, locally uniformly
for \(s\ge0\),
\begin{equation}\label{eq:cf-time-change}
 \frac{\log Q_{\lfloor\kappa Ls\rfloor}}L
 \longrightarrow s
 \quad\text{in probability}.
\end{equation}
Indeed,
\[
 \int_0^1-\log x\,d\nu_G(x)=\frac{\pi^2}{12\log2}
 =\kappa^{-1},
\]
and the two sides of \eqref{eq:cf-roof} differ from \(\log Q_n\)
by at most \(\log2\).  To obtain the claimed uniformity, write
\(\mathscr S_n=\sum_{j<n}-\log x_j\) and
\(\bar r=\kappa^{-1}\).  On every Birkhoff
generic orbit, for any \(C<\infty\) and \(\delta>0\), choose \(n_0\)
so that \(|\mathscr S_n-n\bar r|\le\delta n\) for \(n\ge n_0\).  Then
\[
 \frac1L\max_{n_0\le n\le CL}|\mathscr S_n-n\bar r|\le C\delta,
 \qquad
 \frac1L\max_{n<n_0}|\mathscr S_n-n\bar r|\longrightarrow0.
\]
First let \(L\to\infty\) and then \(\delta\downarrow0\).
Together with the bounded difference above and the harmless floor in
\(\lfloor\kappa Ls\rfloor\), this proves
\eqref{eq:cf-time-change} locally uniformly.  Since Lebesgue and Gauss
measure are equivalent, the same almost-sure statement holds for both.

Equations \eqref{eq:rare-onepoint} and \eqref{eq:cf-time-change},
together with the parity restriction in \eqref{eq:cf-error}, show that
each signed half has intensity
\begin{equation}\label{eq:cf-intensity}
 \frac{\kappa}{2\log2}=\frac6{\pi^2}.
\end{equation}
Thus the deterministic index-time unmarked process has intensity
\(\lambda\), and \eqref{eq:cf-time-change} identifies the candidate
denominator-time intensity.  The rigorous process-level time change is
carried out after marked convergence below.  The uniform factorial
bounds needed in that argument have already been proved above.

\medskip
\noindent\emph{Transfer from Gauss to Lebesgue measure.}
We next make the transfer from Gauss measure to Lebesgue measure
quantitative.  Put
\[
 w=\frac{d\Leb}{d\nu_G}=\log2\,(1+x),\qquad
 \mathscr A_b=\sigma(a_1,\ldots,a_b),
\]
and set
\[
 w_L=\E_{\nu_G}(w\mid\mathscr A_{b_L}).
\]
We choose the version of this conditional expectation which is constant
on each half-open depth-\(b_L\) cylinder and, at every rational endpoint,
assign one of the two adjacent cylinder values.  The endpoints form a
null set for both measures.
A depth-\(b\) continued-fraction cylinder has length
\[
 \frac1{Q_b(Q_b+Q_{b-1})}\le F_{b+1}^{-2}.
\]
Since \(w\) is Lipschitz and this cylinderwise conditional average lies
between the minimum and maximum of \(w\) on the cylinder,
\begin{equation}\label{eq:density-cylinder-approximation}
 \|w-w_L\|_\infty\ll F_{b_L+1}^{-2}=o(1).
\end{equation}
Also \(w_L\) is uniformly bounded and
\(\E_{\nu_G}w_L=\E_{\nu_G}w=1\).

We record an aggregate version of this transfer.  Suppose that
\(X_{\mathbf n}\) is a collection of integrable functions, each
measurable with respect to digits beginning at a position \(d_L\), where
\(d_L-b_L\gg L\), and suppose that
\[
 \sum_{\mathbf n}\E_{\nu_G}|X_{\mathbf n}|=O(1).
\]
Using \(w=w_L+(w-w_L)\), followed by
\eqref{eq:functional-psi-mixing}, gives
\begin{align}
 \sum_{\mathbf n}
 \left|\E_{\rm Leb}X_{\mathbf n}
       -\E_{\nu_G}X_{\mathbf n}\right|
 &\le
 \|w-w_L\|_\infty
 \sum_{\mathbf n}\E_{\nu_G}|X_{\mathbf n}| \notag\\
 &\quad+
 2C\vartheta^{-(d_L-b_L)}
 \E_{\nu_G}|w_L|
 \sum_{\mathbf n}\E_{\nu_G}|X_{\mathbf n}|
 =o(1).
 \label{eq:aggregate-density-transfer}
\end{align}
For deterministic index-time intervals bounded away from zero, the
earliest relevant digit position is \(\gg_\varepsilon L\).  Apply
\eqref{eq:aggregate-density-transfer} specifically to the
digit-replaced tuple functions
\[
 X_{\mathbf n}=\prod_{j=1}^r
 \1_{A_{n_j,j}},\qquad
 A_{n,j}=\{L/a_{n+1}\in J_j\}.
\]
These functions are measurable with respect to the digit future
beginning at the earliest \(a_{n_j+1}\), and their summed
\(L^1(\nu_G)\)-norm is \(O(1)\) by
\eqref{eq:rare-factorial-bound}.  Thus their factorial limits are the
same under Gauss and Lebesgue measure.  The earlier aggregate
\(O(L^{-1})\) replacement of the exact events
\(\{L\theta_{n_j}\in J_j\}\) by \(A_{n_j,j}\) also holds under
Lebesgue measure, because
\[
 \Leb(E)\le\|w\|_\infty\,\nu_G(E)
\]
for every error event \(E\).  This transfers the exact-\(\theta\)
unmarked factorial limits without treating an exact-\(\theta\) event
as a future digit event.  Equivalently, one may first discard the
initial \(2b_L\) positions, whose expected rare count is
\(O(b_L/L)=o(1)\), and retain a gap of \(b_L\).  The lower endpoint is
restored by the direct cell estimate \eqref{eq:endpoint-cell-count}
below.  This proves the same unmarked factorial limits, moment bounds,
and time change for Lebesgue measure.

The following exact one-point calculation is a consistency check and is
not used in the multi-point argument below.  Conditional on a
continued-fraction prefix through \(n\), Lebesgue measure gives the
tail \(x_n\) density
\[
 \frac{1+y_n}{(1+x_ny_n)^2}\,dx_n.
\]
Under \(z=Lx_n/(1+x_ny_n)=L\theta_n\), this becomes exactly
\((1+y_n)L^{-1}\,dz\).  Since
\[
 N(Q_n\alpha-P_n)=(-1)^n\frac{N}{LQ_n}z,
\]
for every nonzero \(h\in\Z\) and compact interval \(J\),
\[
 \left|\int_J e\left(h\frac{Nz}{LQ_n}\right)\,dz\right|
 \le \frac{LQ_n}{2|h|N}.
\]
This tends to zero uniformly on \(Q_n\le N^{1-\varepsilon}\).
The joint independence of several such marks is the content of the
oscillatory argument below.

\emph{The oscillatory marking estimate.}
It remains to show that the torus coordinates in
\eqref{eq:cf-process} become independent Haar marks.  We prove this
through the Fourier coefficients of the factorial measures.
Fix \(r\), compact positive value intervals, prescribed parities, and
compact deterministic index-time intervals contained in
\([\varepsilon,1-\varepsilon]\).  A general labeled \(r\)-th
factorial rectangle is a sum over distinct indices, without an a priori
order.  Partition that sum according to the at most \(r!\) possible
relative orders of the indices, and, in each part, relabel the time
intervals, value intervals, parities, and Fourier frequencies
accordingly.  It is therefore enough to treat one such part.  For
ordered indices
\(n_1<\cdots<n_r\), attach the Fourier factor
\begin{equation}\label{eq:joint-mark-character}
 \exp\left(2\pi\ii N
 \sum_{j=1}^rh_j(Q_{n_j}\alpha-P_{n_j})\right),
 \qquad \mathbf h=(h_1,\ldots,h_r)\in\Z^r.
\end{equation}
We claim that the sum of the corresponding mixed factorial moment
converges to zero whenever \(\mathbf h\ne0\).  When
\(\mathbf h=0\), the unmarked calculation above gives the required
product measure.

Choose
\begin{equation}\label{eq:parameter-hierarchy}
 0<\Delta<\rho/100,\qquad 0<\rho<\varepsilon/100.
\end{equation}
Partition the deterministic time intervals into boxes of diameter at
most \(\Delta\).  For each resulting ordered box configuration, write
\(t_1<\cdots<t_r\) for representative times.  Discard the whole
configuration if \(|t_i-t_j|\le11\rho\) for some \(i\ne j\).
There are finitely many configurations for fixed \(\Delta,\rho\);
the unmarked factorial convergence already proved above shows that
their aggregate mass converges to the corresponding product-Lebesgue
mass.  The latter is \(O_r(\rho)\), since it lies in the union of
\(\binom r2\) diagonal strips of width \(22\rho+2\Delta\).
Thus this deletion costs \(O_r(\rho)+o(1)\).  In every retained
configuration, arbitrary actual times in the chosen boxes differ by
more than
\[
 11\rho-2\Delta>10\rho.
\]
In particular, all ordering and separation restrictions are imposed
by the boxes themselves, rather than by a later tuplewise deletion.

The unmarked factorial bounds also allow us to work on the following
event.  Fix once and for all \(C_*>3\), and set
\begin{equation}\label{eq:good-time-event}
 \mathcal G_{L,\Delta}=
 \left\{
 \sup_{0\le n\le C_*L}
 \left|\frac{\log Q_n}{L}-\frac{n}{\kappa L}\right|
 \le\Delta\right\}.
\end{equation}
The probability of its complement tends to zero by
\eqref{eq:cf-time-change}.  Since the rare counts have bounded moments
of every fixed order, their factorial powers are uniformly integrable.
Indeed, for every fixed \(r\), \eqref{eq:rare-count-moments} and
Cauchy--Schwarz give
\begin{equation}\label{eq:good-event-UI}
 \E\!\left[R_L^r\1_{\mathcal G_{L,\Delta}^c}\right]
 \le(\E R_L^{2r})^{1/2}
    \Pp(\mathcal G_{L,\Delta}^c)^{1/2}=o(1).
\end{equation}
Here \(\E,\Pp\) may refer to either Gauss or Lebesgue measure, because
\eqref{eq:rare-count-moments} transfers as above and
\(d\Leb/d\nu_G\) is bounded.
Therefore inserting or deleting
\(\mathcal G_{L,\Delta}\) changes the factorial moments by \(o(1)\).
On \(\mathcal G_{L,\Delta}\),
\begin{equation}\label{eq:Q-time-box}
 N^{t_j-2\Delta}\le Q_{n_j}\le N^{t_j+2\Delta}.
\end{equation}

Let \(k\) be the largest index for which \(h_k\ne0\).  The time
separation and \eqref{eq:Q-time-box} give
\begin{equation}\label{eq:D-dominance}
 D:=\sum_{j\le k}h_jQ_{n_j},\qquad
 |D|\asymp_{\mathbf h}Q_{n_k}.
\end{equation}
Indeed, for \(j<k\),
\[
 \frac{Q_{n_j}}{Q_{n_k}}
 \le N^{t_j-t_k+4\Delta}
 \le N^{-10\rho+6\Delta}=o(1),
\]
so \(D=h_kQ_{n_k}(1+o(1))\), uniformly over the retained tuples.
The integer numerator terms in \eqref{eq:joint-mark-character}
disappear, so its phase is \(e(ND\alpha)\).

For a block ending at \(n_\ell\), let
\begin{equation}\label{eq:local-denominator-event}
 \mathcal B_{\mathbf n,\ell}
 =\bigcap_{j\le\ell}
 \left\{N^{t_j-2\Delta}\le Q_{n_j}
                         \le N^{t_j+2\Delta}\right\}.
\end{equation}
This event is measurable on the cylinder through \(n_\ell\), and it
contains \(\mathcal G_{L,\Delta}\) for every retained tuple.

\medskip
\noindent\emph{Oscillatory cylinder-sum estimate.}
We isolate the estimate used in both the early and late cases.  Fix
\(1\le\ell\le r\), let
\(\mathscr I_1,\ldots,\mathscr I_\ell\) be the integer index boxes
corresponding to representative times \(t_1<\cdots<t_\ell\), and
assume that \(k\le\ell\) is the largest index for which \(h_k\ne0\).
Thus \(\#\mathscr I_j=O(L)\).  For
\(\mathbf n=(n_1,\ldots,n_\ell)\), put, cylinderwise,
\[
 D_{\mathbf n}=\sum_{j\le k}h_jQ_{n_j},\qquad
 E_{\mathbf n,\ell}
 =\mathcal B_{\mathbf n,\ell}
  \cap\bigcap_{j=1}^{\ell}\{L\theta_{n_j}\in J_j\}.
\]
Then
\begin{equation}\label{eq:oscillatory-cylinder-sum}
 \sum_{\substack{n_j\in\mathscr I_j\\
                  n_1<\cdots<n_\ell}}
 \left|
  \int_0^1
   e\bigl(ND_{\mathbf n}(\alpha)\alpha\bigr)
   \1_{E_{\mathbf n,\ell}}(\alpha)\,d\alpha
 \right|
 \ll_{\mathbf h,r,\mathbf J}
 L^{\ell-1}N^{\,2t_\ell-1-t_k+6\Delta}.
\end{equation}

Here is the complete count.  Decompose the integral for a fixed tuple
into the half-open continued-fraction cylinders \(C\) of depth
\(n_\ell\).  On \(C\), all \(P_{n_j},Q_{n_j}\), \(j\le\ell\), and
hence \(D_{\mathbf n}\), are constant.  Moreover,
\[
 L\theta_{n_j}
 =(-1)^{n_j}LQ_{n_j}(Q_{n_j}\alpha-P_{n_j})
\]
is affine in \(\alpha\).  The intersection of \(C\) with all value
constraints is therefore either empty or a single interval
\(I_{\mathbf n,C}\), and
\(\mathcal B_{\mathbf n,\ell}\) is constant on \(C\).  By the
separation of the representative times and
\(\mathcal B_{\mathbf n,\ell}\),
\[
 |D_{\mathbf n,C}|
 \ge |h_k|Q_{n_k}-\sum_{j<k}|h_j|Q_{n_j}
 \gg_{\mathbf h}N^{t_k-2\Delta}
\]
for all sufficiently large \(N\), uniformly over the retained tuples.
Consequently,
\begin{equation}\label{eq:single-cylinder-oscillation}
 \left|
  \int_{I_{\mathbf n,C}}e(ND_{\mathbf n,C}\alpha)\,d\alpha
 \right|
 \le\frac1{\pi N|D_{\mathbf n,C}|}
 \ll_{\mathbf h}N^{-1-t_k+2\Delta}.
\end{equation}

A finite regular continued-fraction word is determined by its reduced
terminal convergent \(P/Q\), apart from the two usual terminal
expansions.  Indeed, the unique canonical expansion has last digit at
least two, and its only alternative replaces the terminal digit \(a\)
by \(a-1,1\); the Euclidean algorithm gives no other finite expansion.
Hence the number of all finite cylinders whose terminal
denominator is at most \(R\) is at most
\[
 2\sum_{Q\le R}\varphi(Q)+O(1)\ll R^2.
\]
On \(\mathcal B_{\mathbf n,\ell}\) we may take
\(R=N^{t_\ell+2\Delta}\).  Once the deepest word, and hence \(n_\ell\),
is fixed, there are at most \(O(L^{\ell-1})\) choices for
\(n_1,\ldots,n_{\ell-1}\).  Notice that the cylinder count already
includes the summation over \(n_\ell\); no additional factor \(L\) is
present.  Multiplying this count by
\eqref{eq:single-cylinder-oscillation} gives
\[
 L^{\ell-1}N^{2t_\ell+4\Delta}N^{-1-t_k+2\Delta},
\]
which proves \eqref{eq:oscillatory-cylinder-sum}.

Replacing \(\1_{\mathcal G_{L,\Delta}}\) by
\(\1_{\mathcal B_{\mathbf n,\ell}}\) changes the aggregate factorial
expression only on \(\mathcal G_{L,\Delta}^c\), hence by \(o(1)\) in
view of \eqref{eq:good-event-UI}.  We use \(\ell=r\) in the early case
and \(\ell=j_0\) only for the prefix marginal in the late case; no
future denominator condition is passed through a mixing estimate.

\medskip
\noindent\emph{Early and late oscillatory cases.}
Put \(\sigma=(1+t_k)/2\).  Discard every remaining box configuration
for which
\[
 |t_j-\sigma|<4\rho\quad\text{for some }j>k.
\]
The unmarked product factorial measure bounds the prelimit discarded
mass by \(O_r(\rho)+o(1)\).  If every later event has time below
\(\sigma-4\rho\), apply
\eqref{eq:oscillatory-cylinder-sum} with \(\ell=r\); its exponent is
at most \(-8\rho+6\Delta\), so the contribution is \(o(1)\).

It remains to treat the case in which a later event has time above
\(\sigma+4\rho\).  Let \(j_0\ge k\) be the last event below
\(\sigma-4\rho\); include all constraints through \(j_0\) in a prefix
block.  Write
\[
 \mathbf n_-=(n_1,\ldots,n_{j_0}),\qquad
 \mathbf n_+=(n_{j_0+1},\ldots,n_r).
\]
Put \(s=\sigma+\rho\) and use the deterministic depth
\begin{equation}\label{eq:deterministic-depth}
 m_s=\left\lceil\kappa L(s+2\Delta)\right\rceil.
\end{equation}
Define the genuinely prefix-measurable event
\begin{equation}\label{eq:prefix-good-event}
 \mathcal G^-_{L,\Delta}(m_s)=
 \left\{
  \sup_{0\le n\le m_s}
  \left|\frac{\log Q_n}{L}-\frac{n}{\kappa L}\right|
  \le\Delta
 \right\}.
\end{equation}
We have
\(\mathcal G_{L,\Delta}\subset
  \mathcal G^-_{L,\Delta}(m_s)\subset
  \mathcal B_{\mathbf n,j_0}\), and the complement of
\(\mathcal G^-_{L,\Delta}(m_s)\) has probability \(o(1)\).
For the second inclusion, each retained index satisfies
\[
 \left|\frac{n_j}{\kappa L}-t_j\right|\le\Delta;
\]
combining this with the defining bound for
\(\mathcal G^-_{L,\Delta}(m_s)\) gives
\(|L^{-1}\log Q_{n_j}-t_j|\le2\Delta\), exactly the condition in
\(\mathcal B_{\mathbf n,j_0}\).
Here \(m_s<C_*L\) for all large \(N\), by
\(t_k\le1-\varepsilon\) and \eqref{eq:parameter-hierarchy}, so the first
inclusion uses only the range present in \eqref{eq:good-time-event}.
On \(\mathcal G^-_{L,\Delta}(m_s)\),
\[
 Q_{m_s}\ge N^{s+\Delta},\qquad
 m_s>n_{j_0},\qquad
 n_{j_0+1}-m_s\ge c_\rho L.
\]
Indeed,
\[
 \frac{n_{j_0}}{\kappa L}
 \le\sigma-4\rho+\Delta,\qquad
 \frac{m_s}{\kappa L}\ge\sigma+\rho+2\Delta,
\]
while
\[
 \frac{n_{j_0+1}}{\kappa L}
 \ge\sigma+4\rho-\Delta,\qquad
 \frac{m_s}{\kappa L}
 \le\sigma+\rho+2\Delta+O(L^{-1}).
\]
The two index gaps are therefore bounded below by a positive constant
times \(\rho L\).  The lower bound for \(Q_{m_s}\) follows from
\(m_s/(\kappa L)\ge s+2\Delta\) and the defining error \(\Delta\) in
\(\mathcal G^-_{L,\Delta}(m_s)\).

\medskip
\noindent\emph{Late case: freezing the prefix.}
Let \(\mathscr N_-\) denote the retained prefix tuples.  Define first
the shallow-cylinder prefix function
\[
 G^{\mathcal B}_{\mathbf n_-}(\alpha)
 =
 e\bigl(ND_{\mathbf n_-}(\alpha)\alpha\bigr)
 \1_{\mathcal B_{\mathbf n,j_0}}(\alpha)
 \prod_{j\le j_0}
 \1_{\{L\theta_{n_j}(\alpha)\in J_j\}}.
\]
The explicit cylinder estimate
\eqref{eq:oscillatory-cylinder-sum}, now with \(\ell=j_0\), gives
\begin{align}
 \sum_{\mathbf n_-\in\mathscr N_-}
 \left|\E_{\rm Leb}G^{\mathcal B}_{\mathbf n_-}\right|
 &\ll
 L^{j_0-1}N^{2t_{j_0}-1-t_k+6\Delta} \notag\\
 &\le L^{j_0-1}N^{-8\rho+6\Delta}=o(1).
 \label{eq:shallow-prefix-zero}
\end{align}

Next replace the local denominator event by the genuinely
prefix-measurable good event:
\[
 G^-_{\mathbf n_-}(\alpha)
 =
 e\bigl(ND_{\mathbf n_-}(\alpha)\alpha\bigr)
 \1_{\mathcal G^-_{L,\Delta}(m_s)}(\alpha)
 \prod_{j\le j_0}
 \1_{\{L\theta_{n_j}(\alpha)\in J_j\}}.
\]
Since
\(\mathcal G^-_{L,\Delta}(m_s)\subset
  \mathcal B_{\mathbf n,j_0}\),
the two functions agree on
\(\mathcal G^-_{L,\Delta}(m_s)\).  Thus
\begin{align}
 \sum_{\mathbf n_-\in\mathscr N_-}
 \E_{\rm Leb}
 \left|G^-_{\mathbf n_-}
       -G^{\mathcal B}_{\mathbf n_-}\right|
 &\le
 \E_{\rm Leb}\left[
 R_L^{j_0}
 \1_{(\mathcal G^-_{L,\Delta}(m_s))^c}\right] \notag\\
 &\le
 \bigl(\E_{\rm Leb}R_L^{2j_0}\bigr)^{1/2}
 \Leb\bigl((\mathcal G^-_{L,\Delta}(m_s))^c\bigr)^{1/2}
 =o(1).
 \label{eq:prefix-good-replacement}
\end{align}

We now freeze this function on depth-\(m_s\) cylinders.  Let
\(\mathscr C_{m_s}\) be the half-open continued-fraction cylinders of
depth \(m_s\), which partition the irrational points up to a null set.
We keep one fixed half-open endpoint convention and assign all rational
endpoints arbitrarily; every identity in this paragraph is consequently
an almost-everywhere identity, which is all that the integrals require.
For every \(C\in\mathscr C_{m_s}\), choose
an irrational representative \(\alpha_C\in C\), and write
\[
 D_{\mathbf n_-,C}=\sum_{j\le k}h_jQ_{n_j}(C).
\]
The event \(\mathcal G^-_{L,\Delta}(m_s)\) is constant on every such
cylinder.  Define
\begin{equation}\label{eq:frozen-prefix-definition}
 F_{\mathbf n_-}(\alpha)
 =
 \sum_{C\in\mathscr C_{m_s}}
 \1_C(\alpha)
 \1_{\{C\subset\mathcal G^-_{L,\Delta}(m_s)\}}
 e\bigl(ND_{\mathbf n_-,C}\alpha_C\bigr)
 \prod_{j\le j_0}
 \1_{\{L\theta_{n_j}(\alpha_C)\in J_j\}}.
\end{equation}
This function is bounded and
\(\sigma(a_1,\ldots,a_{m_s})\)-measurable.

A depth-\(m_s\) cylinder has length at most \(Q_{m_s}^{-2}\).
On a cylinder retained by
\(\mathcal G^-_{L,\Delta}(m_s)\), the denominator bounds above imply
\begin{equation}\label{eq:phase-freeze}
 \sup_{\alpha\in C}
 \left|
 e(ND_{\mathbf n_-,C}\alpha)
 -e(ND_{\mathbf n_-,C}\alpha_C)
 \right|
 \ll_{\mathbf h}
 \frac{N^{1+t_k+2\Delta}}{N^{2s+2\Delta}}
 =O_{\mathbf h}(N^{-2\rho})
 =:\delta_N.
\end{equation}
Similarly, since \(t_j\le t_{j_0}\le\sigma-4\rho\),
\begin{align}
 \sup_{\alpha\in C}
 \left|L\theta_{n_j}(\alpha)
       -L\theta_{n_j}(\alpha_C)\right|
 &\le\frac{LQ_{n_j}^2}{Q_{m_s}^2} \notag\\
 &\ll
 L N^{2(\sigma-4\rho+2\Delta)
       -2(\sigma+\rho+\Delta)}
 \le L N^{-10\rho+2\Delta}
 =:\eta_N
 \qquad(j\le j_0).
 \label{eq:value-freeze}
\end{align}

For an interval \(J\), write \(J^{(\eta)}\) for its closed
\(\eta\)-neighborhood.  The two freezing estimates give the pointwise
bound
\begin{align}
 |G^-_{\mathbf n_-}-F_{\mathbf n_-}|
 &\ll
 \delta_N
 \prod_{j\le j_0}
 \1_{\{L\theta_{n_j}\in J_j^{(\eta_N)}\}} \notag\\
 &\quad+
 \sum_{i\le j_0}
 \1_{\{\operatorname{dist}
          (L\theta_{n_i},\partial J_i)\le\eta_N\}}
 \prod_{\substack{j\le j_0\\j\ne i}}
 \1_{\{L\theta_{n_j}\in J_j^{(\eta_N)}\}}.
 \label{eq:freezing-pointwise-error}
\end{align}
For every fixed \(z>0\), uniformly for
\(0<\eta\le\min(1,z/2)\),
\begin{equation}\label{eq:value-boundary-strip}
 \nu_G\{|L\theta_n-z|\le\eta\}
 \ll_z\frac{\eta}{L}+\frac1{L^2}.
\end{equation}
Indeed, on this event \(L\theta_n\) remains in a fixed compact subset
of \((0,\infty)\), and \eqref{eq:digit-replacement} places it inside
\[
 \left\{\left|\frac{L}{a_{n+1}}-z\right|
             \le\eta+\frac{C_z}{L}\right\}.
\]
Direct summation of the one-digit probabilities proves
\eqref{eq:value-boundary-strip}.

For intersections with the other rare constraints, enlarge those
constraints to their corresponding one-digit events and use
\eqref{eq:digit-quasi-bernoulli}.  There are \(O(L^{j_0})\) prefix
tuples, so the total contribution of any one boundary coordinate is
\[
 O\left(
 L^{j_0}
 \left(\frac{\eta_N}{L}+\frac1{L^2}\right)
 L^{-(j_0-1)}
 \right)
 =O(\eta_N+L^{-1}).
\]
The first term in \eqref{eq:freezing-pointwise-error} contributes
\(O(\delta_N)\) by \eqref{eq:rare-factorial-bound}.  The densities
of Gauss and Lebesgue measure relative to one another are bounded, so
the estimates hold under either measure.  Consequently,
\begin{equation}\label{eq:aggregate-freezing-error}
 \sum_{\mathbf n_-\in\mathscr N_-}
 \E_{\rm Leb}|G^-_{\mathbf n_-}-F_{\mathbf n_-}|
 \ll\delta_N+\eta_N+L^{-1}=o(1).
\end{equation}
The same calculation with all \(r\) rare factors present gives
\begin{equation}\label{eq:full-aggregate-freezing-error}
 \sum_{\mathbf n_-,\mathbf n_+}
 \E_{\rm Leb}\left(
 |G^-_{\mathbf n_-}-F_{\mathbf n_-}|
 \prod_{j>j_0}\1_{\{L\theta_{n_j}\in J_j\}}\right)
 =o(1).
\end{equation}
More precisely, after enlarging every exact rare event to its
one-digit event, the phase term in this sum is
\[
 O(\delta_N L^rL^{-r})=O(\delta_N),
\]
whereas a term with one prefix boundary strip is
\[
 O\left(
 L^r\left(\frac{\eta_N}{L}+\frac1{L^2}\right)
 L^{-(r-1)}\right)
 =O(\eta_N+L^{-1}).
\]
There are only \(j_0\le r\) possible boundary coordinates.  This proves
\eqref{eq:full-aggregate-freezing-error} without multiplying a merely
prefix-level error by the number of future tuples.
Combining
\eqref{eq:shallow-prefix-zero},
\eqref{eq:prefix-good-replacement}, and
\eqref{eq:aggregate-freezing-error} proves the required absolute
summability, because
\[
 \sum_{\mathbf n_-}|\E_{\rm Leb}F_{\mathbf n_-}|
 \le
 \sum_{\mathbf n_-}|\E_{\rm Leb}G^{\mathcal B}_{\mathbf n_-}|
 +\sum_{\mathbf n_-}\E_{\rm Leb}
       |G^-_{\mathbf n_-}-G^{\mathcal B}_{\mathbf n_-}|
 +\sum_{\mathbf n_-}\E_{\rm Leb}
       |F_{\mathbf n_-}-G^-_{\mathbf n_-}|.
\]
Thus
\begin{equation}\label{eq:prefix-zero}
 \sum_{\mathbf n_-\in\mathscr N_-}
 \left|\E_{\rm Leb}F_{\mathbf n_-}\right|=o(1).
\end{equation}

The frozen factor is itself dominated by enlarged one-digit rare
events.  In fact, if \(F_{\mathbf n_-}(\alpha)\ne0\) and
\(\alpha\in C\), then
\eqref{eq:value-freeze} gives
\[
 L\theta_{n_j}(\alpha)\in J_j^{(\eta_N)}
 \qquad(j\le j_0).
\]
Applying \eqref{eq:digit-replacement} on this fixed enlarged compact
window, with a constant uniform in \(j\le j_0\), yields
\begin{equation}\label{eq:frozen-prefix-digit-domination}
 |F_{\mathbf n_-}(\alpha)|
 \le
 \prod_{j\le j_0}
 \1_{\{L/a_{n_j+1}\in
              J_j^{(\eta_N+C_{\mathbf J}/L)}\}}.
\end{equation}

\medskip
\noindent\emph{Late case: prefix--future mixing.}
All later mark frequencies are zero.  By
\eqref{eq:digit-replacement}, their exact value constraints can be
replaced, with aggregate error \(o(1)\), by the one-digit constraints
\(L/a_{n+1}\in J\).  To see that this remains true after multiplication
by \(F_{\mathbf n_-}\), use
\eqref{eq:frozen-prefix-digit-domination}.  If a single future
indicator changes, its digit lies in an endpoint strip of probability
\(O(L^{-2})\); the frozen prefix and the remaining future coordinates
are bounded by \(r-1\) one-digit rare events of probability
\(O(L^{-1})\).  Hence \eqref{eq:digit-quasi-bernoulli} gives the full
aggregate bound
\[
 O\bigl(L^rL^{-2}L^{-(r-1)}\bigr)=O(L^{-1}).
\]
Let
\(H_{\mathbf n_+}\) denote the product of these future one-digit
indicators, and let \(\mathscr N_+\) be the retained future tuples.
Then \(H_{\mathbf n_+}\) is measurable with respect to digits beginning
at \(a_{n_{j_0+1}+1}\), at distance at least \(c_\rho L\) from the
prefix \(\sigma\)-field.
For the fixed separated collection of index boxes now under
consideration, every retained prefix tuple precedes every retained
future tuple, and the boxes themselves enforce all ordering and
separation conditions.  Hence the retained full tuples are exactly
the Cartesian product \(\mathscr N_-\times\mathscr N_+\), which
justifies the product double sums below.

We also replace the global event \(\mathcal G_{L,\Delta}\) by
\(\mathcal G^-_{L,\Delta}(m_s)\).  Since the former is contained in
the latter, note explicitly that on any fixed digit window
\(L/a_{n+1}\in K_0\Subset(0,\infty)\),
\[
 0\le \frac{L}{a_{n+1}}-L\theta_n
 =\frac{L(x_{n+1}+y_n)}
        {a_{n+1}(a_{n+1}+x_{n+1}+y_n)}
 \le\frac{2L}{a_{n+1}^2}=O_{K_0}(L^{-1}).
\]
Thus all digit-replaced events occurring below are contained, for
large \(L\), in a fixed enlargement of the corresponding exact
\(\theta\)-windows.  Take \(R_L\) here on a fixed slightly enlarged compact
value window containing every exact prefix window, every enlargement
in \eqref{eq:frozen-prefix-digit-domination}, and every
digit-replaced future window.  It therefore dominates every summand
that occurs in this replacement, and the aggregate difference is at
most
\[
 \E\left[R_L^r\1_{\mathcal G_{L,\Delta}^c}\right]=o(1)
\]
by \eqref{eq:good-event-UI}.  Together with
\eqref{eq:full-aggregate-freezing-error}, this shows that the late part
of the original factorial Fourier coefficient is
\begin{equation}\label{eq:late-frozen-sum}
 \sum_{\mathbf n_-\in\mathscr N_-}
 \sum_{\mathbf n_+\in\mathscr N_+}
 \E_{\rm Leb}\bigl(F_{\mathbf n_-}H_{\mathbf n_+}\bigr)+o(1).
\end{equation}

Recall \(w=d\Leb/d\nu_G\) and
\(w_L=\E_{\nu_G}(w\mid a_1,\ldots,a_{b_L})\).  Since \(b_L<m_s\),
\(w_LF_{\mathbf n_-}\) is prefix measurable.  The estimates just
proved, together with \eqref{eq:rare-factorial-bound}, give
\[
 \sum_{\mathbf n_-}\E_{\nu_G}|w_LF_{\mathbf n_-}|=O(1),
 \qquad
 \sum_{\mathbf n_+}\E_{\nu_G}H_{\mathbf n_+}=O(1).
\]
The functional mixing inequality
\eqref{eq:functional-psi-mixing} now yields the aggregate covariance
bound
\begin{align}
 &\sum_{\mathbf n_-,\mathbf n_+}
 \left|
 \E_{\nu_G}(w_LF_{\mathbf n_-}H_{\mathbf n_+})
 -\E_{\nu_G}(w_LF_{\mathbf n_-})
  \E_{\nu_G}H_{\mathbf n_+}
 \right| \notag\\
 &\qquad\le
 2C\vartheta^{-c_\rho L}
 \left(\sum_{\mathbf n_-}
       \E_{\nu_G}|w_LF_{\mathbf n_-}|\right)
 \left(\sum_{\mathbf n_+}
       \E_{\nu_G}H_{\mathbf n_+}\right)
 =o(1).
 \label{eq:aggregate-late-mixing}
\end{align}

Applying the same inequality to \(|F_{\mathbf n_-}|\) and
\(H_{\mathbf n_+}\) shows that
\[
 \sum_{\mathbf n_-,\mathbf n_+}
 \E_{\nu_G}|F_{\mathbf n_-}|H_{\mathbf n_+}=O(1).
\]
Thus replacing \(w\) by \(w_L\) in
\eqref{eq:late-frozen-sum} costs at most
\[
 \|w-w_L\|_\infty
 \sum_{\mathbf n_-,\mathbf n_+}
 \E_{\nu_G}|F_{\mathbf n_-}|H_{\mathbf n_+}=o(1).
\]
Moreover,
\begin{equation}\label{eq:prefix-density-error}
 \sum_{\mathbf n_-}
 \left|
 \E_{\nu_G}(w_LF_{\mathbf n_-})
 -\E_{\rm Leb}F_{\mathbf n_-}
 \right|
 \le
 \|w-w_L\|_\infty
 \sum_{\mathbf n_-}\E_{\nu_G}|F_{\mathbf n_-}|
 =o(1).
\end{equation}
For completeness,
\eqref{eq:aggregate-density-transfer}, applied to the future tuple
indicators, also gives
\begin{equation}\label{eq:future-density-error}
 \sum_{\mathbf n_+}
 \left|
 \E_{\rm Leb}H_{\mathbf n_+}
 -\E_{\nu_G}H_{\mathbf n_+}
 \right|=o(1).
\end{equation}

Combining
\eqref{eq:late-frozen-sum},
\eqref{eq:aggregate-late-mixing},
\eqref{eq:prefix-density-error}, and
\eqref{eq:prefix-zero}, the absolute value of the late contribution is
at most
\[
 \left(\sum_{\mathbf n_-}
       |\E_{\rm Leb}F_{\mathbf n_-}|\right)
 \left(\sum_{\mathbf n_+}
       \E_{\nu_G}H_{\mathbf n_+}\right)+o(1)
 =o(1).
\]
Finally let \(N\to\infty\), then \(\Delta\downarrow0\), and then
\(\rho\downarrow0\).  We have proved that every nonzero joint Fourier
coefficient of the marked factorial measures vanishes.

\medskip
\noindent\emph{Signed extension and the point-process criterion.}
We have so far written the argument for positive value intervals and
prescribed parity.  This covers the whole signed state space.  Indeed,
by \eqref{eq:cf-error},
\[
 (-1)^nL\theta_n>0
 \quad\Longleftrightarrow\quad n\ \hbox{is even}.
\]
Thus an interval \(J\Subset(0,\infty)\) selects even indices, whereas
an interval \(J\Subset(-\infty,0)\) selects odd indices after reflection
to the positive interval \(-J\).  Reflection preserves Lebesgue length,
so both signed halves have the same intensity.  Every compact subset of
\(\R\setminus\{0\}\) has positive distance from zero and is the union
of its positive and negative parts.  Finite disjoint decompositions
therefore reduce every relatively compact signed rectangle to the two
cases just proved.  No rectangle accumulating at zero is needed,
because such a set is not relatively compact in
\(\R\setminus\{0\}\).

Trigonometric polynomials are uniformly dense in the continuous
functions on \(\T^r\).  The zero Fourier coefficient is the unmarked
product factorial measure, and all other coefficients vanish.  Hence,
the factorial measures converge weakly in the torus variables.  For a
Haar-continuity rectangle, sandwich its torus indicator between continuous
functions whose Haar integrals differ by an arbitrarily small amount; the
uniform unmarked factorial bound controls the corresponding aggregate
error.  Therefore,
on \([\varepsilon,1-\varepsilon]\), all limiting factorial measures
are
\[
 \lambda^r\prod_{j=1}^r dt_j\,d\xi_j\,dm(u_j).
\]
We now spell out the passage from factorial measures to point
processes.  Let
\[
 \widetilde{\mathcal Q}_N
 =
 \sum_{n\ge0}
 \delta_{\left(n/(\kappa L),\;(-1)^nL\theta_n,\;
                    (-1)^nN\theta_n/Q_n\bmod1\right)}
\]
be the process in deterministic index time.  On
\[
 E_\varepsilon
 =[\varepsilon,1-\varepsilon]
  \times(\R\setminus\{0\})\times\T,
\]
let \(\mathscr R_\varepsilon\) be the countable semiring of relatively
compact half-open rectangles whose Euclidean endpoints are rational,
whose value-coordinate closures avoid zero, and whose torus endpoints
belong to a fixed countable dense subset of \(\T\).  A torus arc that
wraps around zero is split into two arcs.  This is a dissecting,
convergence-determining semiring, and every member is a continuity set
for the diffuse measure
\[
 \Lambda=\lambda\,dt\,d\xi\,dm.
\]

Take pairwise disjoint
\(B_1,\ldots,B_q\in\mathscr R_\varepsilon\), and write
\(X_{N,i}=\widetilde{\mathcal Q}_N(B_i)\).  The
factorial-measure calculation above
gives, for every \(s_1,\ldots,s_q\ge0\),
\begin{equation}\label{eq:joint-count-factorial-moments}
 \E\prod_{i=1}^q(X_{N,i})_{s_i}
 \longrightarrow
 \prod_{i=1}^q\Lambda(B_i)^{s_i}.
\end{equation}
The right-hand side is the joint factorial moment of independent
Poisson variables \(Z_i\) with means \(\Lambda(B_i)\).

Let \(T_N=\sum_iX_{N,i}\).  After reflecting every negative value
window to its positive magnitude, \(T_N\) is bounded by a finite union
of rare counts of the form used in \eqref{eq:rare-factorial-bound}.
Consequently, for every fixed \(m\),
\[
 \sup_N\E(T_N)_m<\infty,\qquad
 \sup_N\E T_N^m<\infty,
\]
where the second bound follows from
\eqref{eq:ordinary-factorial-identity}.  In particular, the bound with
one order larger makes every polynomial of a fixed degree in
\((X_{N,1},\ldots,X_{N,q})\) uniformly integrable.  Every ordinary
joint moment is a finite linear combination of the joint factorial
moments in \eqref{eq:joint-count-factorial-moments}; hence all ordinary
joint moments converge to those of \((Z_1,\ldots,Z_q)\).  That law is
moment determinate because its moment-generating function
\[
 \exp\left(\sum_{i=1}^q
       \Lambda(B_i)(e^{t_i}-1)\right)
\]
is finite in a neighborhood of the origin.  The multivariate method of
moments therefore yields
\begin{equation}\label{eq:joint-poisson-counts}
 (X_{N,1},\ldots,X_{N,q})
 \Longrightarrow (Z_1,\ldots,Z_q).
\end{equation}

For every finite union \(U\) of members of
\(\mathscr R_\varepsilon\), we have \(\Lambda(\partial U)=0\), and the
first factorial moment and
\eqref{eq:joint-poisson-counts} give
\[
 \E\widetilde{\mathcal Q}_N(U)\longrightarrow\Lambda(U),\qquad
 \Pp\{\widetilde{\mathcal Q}_N(U)=0\}
 \longrightarrow e^{-\Lambda(U)},
\]
respectively.  Kallenberg's mean-and-void criterion states that these
two limits on the finite unions of a countable dissecting semiring of
continuity sets imply convergence to the Poisson random measure
\cite[Theorem~4.7]{Kallenberg}.  Applied to
\(\mathscr R_\varepsilon\) on this locally compact second-countable
space, it therefore gives vague convergence to the Poisson process with
intensity \(\Lambda\).  Notice also that the mean
bound supplies local tightness directly: every compact set is covered
by such a finite union \(U\), and
\(\sup_N\E\widetilde{\mathcal Q}_N(U)<\infty\), so Markov's inequality
controls the probability of more than \(M\) points uniformly as
\(M\to\infty\).  Hence the marked point process converges in
deterministic index time.

\medskip
\noindent\emph{Time change and endpoint restoration.}
We give the time-change estimate at the level of test functions.  Let
\(f\) be continuous with compact support whose time projection lies in
\([\varepsilon,1-\varepsilon]\), and let
\[
 \omega_f(\delta)=
 \sup_{\substack{|t-t'|\le\delta\\ \xi\in\R\setminus\{0\},\,u\in\T}}
 |f(t,\xi,u)-f(t',\xi,u)|.
\]
Extend \(f\) by zero outside a slightly larger compact time interval;
then \(\omega_f(\delta)\to0\) as \(\delta\downarrow0\).
Because the value projection of \(\operatorname{supp}f\) is compact in
\(\R\setminus\{0\}\), choose a compact interval
\(K_f\Subset(0,\infty)\) containing
\[
 \bigl\{|\xi|:(t,\xi,u)\in\operatorname{supp}f\bigr\}.
\]
Let \(R_L^{(f)}\) count all rare convergents with
\(0\le n\le C_*L\) and magnitude \(L\theta_n\in K_f\).  This
contains every index-time point relevant to \(f\).  It also contains
every denominator-time point relevant to \(f\), because
\(Q_n\ge F_{n+1}\) and \(Q_n<N\) imply \(n<C_*L\) for all sufficiently
large \(N\).  By
\eqref{eq:rare-count-moments},
\[
 \sup_L\E R_L^{(f)}+\sup_L\E(R_L^{(f)})^2<\infty.
\]

On \(\mathcal G_{L,\Delta}\), with
\(\Delta<\varepsilon/2\), every relevant convergent satisfies
\[
 \left|\frac{n}{\kappa L}-\frac{\log Q_n}{L}\right|\le\Delta.
\]
Moreover, a point with either time coordinate in
\(\operatorname{supp}f\) has the other coordinate in
\([\varepsilon/2,1-\varepsilon/2]\).  In particular \(Q_n<N\), so the
cutoff in \eqref{eq:cf-process} is present on both sides.  Using the
continued-fraction identification \eqref{eq:cf-process}, we obtain
\begin{align}
 \left|
 \langle\widetilde{\mathcal Q}_N,f\rangle
 -\langle\mathcal Q_N,f\rangle
 \right|
 &\le \omega_f(\Delta)R_L^{(f)}
 \qquad\text{on }\mathcal G_{L,\Delta}.
 \label{eq:test-function-time-change}
\end{align}
On the complement, the same difference is at most
\(2\|f\|_\infty R_L^{(f)}\).  Hence
\begin{align*}
 \E\left|
 \langle\widetilde{\mathcal Q}_N,f\rangle
 -\langle\mathcal Q_N,f\rangle
 \right|
 &\le
 \omega_f(\Delta)\E R_L^{(f)}
 +2\|f\|_\infty
   \E\left[R_L^{(f)}\1_{\mathcal G_{L,\Delta}^c}\right]\\
 &\le O_f(\omega_f(\Delta))+o_{N\to\infty}(1),
\end{align*}
where Cauchy--Schwarz, the second-moment bound, and
\(\Pp(\mathcal G_{L,\Delta}^c)=o(1)\) were used in the last line.
First let \(N\to\infty\), and then let \(\Delta\downarrow0\).
The converging-together argument for a countable
convergence-determining family of such \(f\)'s transfers the marked
process from index time to denominator time on every fixed interior
layer.  The endpoint layers are restored next by the direct cell
bound.

It remains only to restore the two endpoint layers.  Here no time-change
argument is needed.  Directly subdividing \(\alpha\) into the reduced
rational cells about \(q/p\) gives, for every set of denominators
\(I\subset[1,N)\),
\begin{equation}\label{eq:endpoint-cell-count}
 \E\#\left\{p\in I:(p,q_p)=1,
                  |Lp\delta_p|\le B\right\}
 \le \frac{2B}{L}\sum_{p\in I}\frac{\varphi(p)}{p^2}
      +O_B(L^{-1}).
\end{equation}
Indeed, for each of the \(\varphi(p)\) reduced numerators the relevant
\(\alpha\)-interval has length at most \(2B/(Lp^2)\).  For \(p>1\)
these reduced cells are interior; the \(O_B(L^{-1})\) error covers the
two endpoint cells for \(p=1\).  We also have
\[
 \sum_{p\le x}\frac{\varphi(p)}{p^2}
 =\frac6{\pi^2}\log x+O(1),
\]
as follows by inserting
\[
 \frac{\varphi(p)}p=\sum_{d\mid p}\frac{\mu(d)}d
\]
and interchanging the sums.  The resulting inner sum uses
\(\sum_{m\le y}m^{-1}=\log y+O(1)\), while the leading coefficient is
\(\sum_{d\ge1}\mu(d)/d^2=1/\zeta(2)\).
Therefore the expected number of points in either discarded time layer
\([0,\varepsilon]\) or \([1-\varepsilon,1]\), on a fixed
\(\xi\)-window, is \(O_B(\varepsilon)+o(1)\).  More explicitly, if
\(f\in C_c([0,1]\times(\R\setminus\{0\})\times\T)\), its value
projection is contained in \(\{|\xi|\le B_f\}\) for some \(B_f\).
Choose a continuous function
\(\chi_\varepsilon:[0,1]\to[0,1]\) which is zero on
\([0,\varepsilon]\cup[1-\varepsilon,1]\) and one on
\([2\varepsilon,1-2\varepsilon]\).  Then
\(f(t,\xi,u)\chi_\varepsilon(t)\) is a compactly supported continuous
test function on an interior layer, and
\[
 \E\left|
 \langle\mathcal Q_N,f\rangle
 -\langle\mathcal Q_N,
          f\chi_\varepsilon\rangle
 \right|
 \le \|f\|_\infty\bigl(O_{B_f}(\varepsilon)+o(1)\bigr).
\]
The limiting Poisson process satisfies the analogous bound
\(O_f(\varepsilon)\) by its intensity measure.  Markov's inequality
and the converging-together argument therefore restore the two layers.
Letting \(\varepsilon\downarrow0\) proves
\eqref{eq:PPP-intensity-time}.  Projection in the \(t\)-coordinate is
proper because \([0,1]\) is compact, so the point-process mapping theorem
proves \eqref{eq:PPP-intensity}.
\end{proof}
 
\subsection{Finite Poisson shots}

For \(A>0\), define
\begin{equation}\label{eq:XNA}
 X_{N,A}=
 \sum_{\substack{p\le N,\ (p,q_p)=1\\|Lp\delta_p|\le A}}
 \frac{V(N\delta_p)}{Lp\delta_p}.
\end{equation}
The \(p=N\) term may be deleted: the measure of its support is
\(O_A((LN)^{-1})\).

Continue to write \(\mathcal P=\{(\xi_j,u_j)\}\) for the Poisson process with
intensity \eqref{eq:PPP-intensity}, and put
\begin{equation}\label{eq:XA}
 X_A=\sum_{|\xi_j|\le A}\frac{V(u_j)}{\xi_j}.
\end{equation}
The sum is almost surely finite and no point has \(\xi_j=0\).

\begin{lemma}[Finite-shot convergence]\label{lem:finite-shot}
For every fixed \(A>0\),
\[
 X_{N,A}\Longrightarrow X_A.
\]
\end{lemma}

\begin{proof}
For \(0<r<A\), restrict both point processes to
\(r\le|\xi|\le A\).  The shot functional is continuous at every simple
configuration with no point on the boundary, an event of limiting
probability one.  Lemma~\ref{lem:marked-poisson} and the continuous
mapping theorem give convergence of the restricted shots.

It remains to let \(r\downarrow0\).  For the prelimit, a union bound and
subdivision into reduced cells give
\[
 \Leb\{\mathcal P_N((-r,r)\times\T)>0\}
 \le
 \frac{2r}{L}\sum_{p\le N}\frac{\varphi(p)}{p^2}
 \ll r.
\]
The limiting probability is at most
\(2r/\zeta(2)\).  With probability \(1-O(r)\), deleting
\(|\xi|<r\) changes neither shot.  Let \(r\downarrow0\).
\end{proof}

\subsection{The Cauchy exponent}

\begin{lemma}[Poisson integral]\label{lem:poisson-Cauchy}
As \(A\to\infty\),
\[
 X_A\Longrightarrow X,\qquad
 \E e^{\ii tX}=\exp\left(-\frac{|t|}{2\pi}\right).
\]
\end{lemma}

\begin{proof}
The exponential formula for a Poisson process and symmetry in \(\xi\)
give
\begin{align*}
 \log\E e^{\ii tX_A}
 &=
 \frac1{\zeta(2)}
 \int_{-A}^{A}\int_0^1
 \left(e^{\ii tV(u)/\xi}-1\right)\dd u\dd\xi\\
 &=
 \frac2{\zeta(2)}
 \int_0^A\int_0^1
 \left(\cos(tV(u)/\xi)-1\right)\dd u\dd\xi.
\end{align*}
For \(a\ge0\),
\[
 \int_0^\infty\left(\cos(a/\xi)-1\right)\dd\xi
 =
 a\int_0^\infty\frac{\cos y-1}{y^2}\dd y
 =-\frac{\pi a}{2}.
\]
Moreover,
\[
 0\le1-\cos(tV(u)/\xi)
 \le\min\!\left(2,\frac{t^2V(u)^2}{2\xi^2}\right)
 \le\min(2,C_t\xi^{-2}),
\]
which is integrable for \(\xi>0\).  Dominated convergence in the
nonpositive symmetric integrand therefore yields
\[
 \lim_{A\to\infty}\log\E e^{\ii tX_A}
 =
 -\frac{\pi|t|}{\zeta(2)}\int_0^1V(u)\dd u.
\]
Since
\[
 \int_0^1V(u)\dd u=\frac1{12},\qquad
 \zeta(2)=\frac{\pi^2}{6},
\]
the exponent is \(-|t|/(2\pi)\).  L\'evy's continuity theorem
completes the proof.
\end{proof}

\subsection{Completion of the proof}

\begin{proof}[Proof of Theorem~\ref{thm:main}]
By Proposition~\ref{prop:natural-cutoff},
\[
 \frac{S_N-Y_N}{L}\longrightarrow0
 \quad\text{in }L^2,\ \text{hence in probability}.
\]
For fixed \(A\), the part of \(Y_N/L\) with
\(|Lp\delta_p|\le A\) is \(X_{N,A}\).  Proposition~\ref{prop:minor}
states precisely that the complementary part goes to zero in probability
when first \(N\to\infty\) and then \(A\to\infty\).
Lemma~\ref{lem:finite-shot} gives
\(X_{N,A}\Rightarrow X_A\) for fixed \(A\), and
Lemma~\ref{lem:poisson-Cauchy} gives \(X_A\Rightarrow X\).
Moreover, for every \(\eta>0\), the union bound gives
\begin{align*}
 \Pp\left\{\left|\frac{S_N}{L}-X_{N,A}\right|>\eta\right\}
 &\le
 \Pp\left\{\frac{|S_N-Y_N|}{L}>\frac\eta2\right\}\\
 &\quad+
 \Pp\left\{\left|\frac{Y_N}{L}-X_{N,A}\right|>\frac\eta2\right\}.
\end{align*}
Consequently,
\[
 \lim_{A\to\infty}\limsup_{N\to\infty}
 \Pp\left\{\left|\frac{S_N}{L}-X_{N,A}\right|>\eta\right\}=0.
\]
The converging-together theorem \cite[Theorem~3.2]{Billingsley} now yields
\[
 \frac{S_N}{\log N}\Longrightarrow X.
\]
The Cauchy law with characteristic function
\(\exp(-|t|/(2\pi))\) has continuous distribution function
\[
 g(c)=\frac12+\frac1\pi\arctan(2\pi c).
\]
Weak convergence therefore gives convergence of the distribution
functions at every \(c\in\R\).  The function \(g\) is nondecreasing,
with limits \(0\) and \(1\) at \(-\infty\) and \(+\infty\),
respectively.
\end{proof}

\section{Two elementary consistency checks}

For completeness, the finite-\(N\) law has no atoms.  Away from the
finitely many points \(a/k\), \(k\le N\),
\[
 S_N(\alpha)=\frac N2-\frac{N(N+1)}2\alpha+
 \sum_{k=1}^N\floor{k\alpha}.
\]
It is therefore affine with nonzero slope on every complementary
interval.  Also, away from those points,
\[
 S_N(1-\alpha)=-S_N(\alpha).
\]
Thus every finite-\(N\) law is exactly symmetric and
\(\Leb\{S_N\le0\}=1/2\), in agreement with the limiting Cauchy law.

The second moment is not uniformly integrable on the logarithmic scale.
Fourier orthogonality gives
\[
 \int_0^1\psi(k\alpha)\psi(\ell\alpha)\dd\alpha
 =\frac{(k,\ell)^2}{12k\ell},
\]
so positivity of all covariance terms and the diagonal terms give
\[
 \E\left(\frac{S_N}{\log N}\right)^2
 =\frac{\Var S_N}{(\log N)^2}
 \ge\frac{N}{12(\log N)^2}\longrightarrow\infty.
\]
Thus \(\{(S_N/\log N)^2\}_N\) is not uniformly integrable.  This
divergence of second moments does not contradict the weak convergence
in Theorem~\ref{thm:main}.

\begin{thebibliography}{99}
\raggedright

\bibitem{Birkhoff}
G.~D. Birkhoff,
Proof of the ergodic theorem,
\emph{Proc. Natl. Acad. Sci. USA} \textbf{17} (1931), 656--660,
doi:10.1073/pnas.17.12.656.

\bibitem{Billingsley}
P.~Billingsley,
\emph{Convergence of Probability Measures}, second edition,
Wiley, New York, 1999.

\bibitem{ChanKumchev}
T.~H. Chan and A.~V. Kumchev,
On sums of Ramanujan sums,
\emph{Acta Arith.} \textbf{152} (2012), 1--10,
doi:10.4064/aa152-1-1.

\bibitem{Erdos1964}
P.~Erd\H{o}s,
Problems and results on diophantine approximations,
\emph{Compositio Math.} \textbf{16} (1964), 52--65,
\url{https://www.numdam.org/item/CM_1964__16__52_0/}.

\bibitem{GhoshKirsebomRoy}
A.~Ghosh, M.~S. Kirsebom, and P.~Roy,
Continued fractions, the Chen--Stein method and extreme value theory,
\emph{Ergodic Theory Dynam. Systems} \textbf{41} (2021), 461--470,
doi:10.1017/etds.2019.64.

\bibitem{Kallenberg}
O.~Kallenberg,
\emph{Random Measures}, third edition,
Akademie-Verlag, Berlin, 1983.

\bibitem{Kesten1960}
H.~Kesten,
Uniform distribution mod \(1\),
\emph{Ann.\ of Math. (2)} \textbf{71} (1960), 445--471.

\bibitem{Kesten1962}
H.~Kesten,
Uniform distribution mod \(1\). II,
\emph{Acta Arith.} \textbf{7} (1961/1962), 355--380.

\bibitem{Khinchin}
A.~Ya. Khinchin,
\emph{Continued Fractions},
University of Chicago Press, Chicago, 1964.

\end{thebibliography}

\end{document}
